% MSc dissertation example file, February 2022
%
% Leave one of the documentclass lines uncommented to match your degree.
% You may remove the logo option if it causes problems.
% Do not change any other options.
% \documentclass[logo,msc,adi]{infthesis}     % Adv Design Inf
% \documentclass[logo,msc,ai]{infthesis}      % AI
% \documentclass[logo,msc,cogsci]{infthesis}  % Cognitive Sci
% \documentclass[logo,msc,cs]{infthesis}      % Computer Sci
% \documentclass[logo,msc,cyber]{infthesis}   % Cyber Sec
% \documentclass[logo,msc,datasci]{infthesis} % Data Sci
% \documentclass[logo,msc,di]{infthesis}      % Design Inf
% \documentclass[logo,msc,dsti]{infthesis}    % Data Sci TI
% \documentclass[logo,msc,inf]{infthesis}     % Informatics
\documentclass[logo,msc,ai]{infthesis}           % degree unspecified, do not change except to add your degree
%%%%%%%%%%%%%%%%%%%%%%%%
% Understand any problems and seek approval before assuming it's ok to remove ugcheck.
% \usepackage{msccheck}
\usepackage{msccheck} 

% Include any packages you need below, but don't include any that change the page
% layout or style of the dissertation. By including the ugcheck package above,
% you should catch most accidental changes of page layout though.

\usepackage{microtype} % recommended, but you can remove if it causes problems
\usepackage{xcolor} %remove later this is used for helping edit and draft
\usepackage{caption}
\usepackage{subcaption} %subcaptions for plots and figures next to each other. 
\usepackage{amsmath}  %Check we are allowed
\usepackage{amssymb} %Check we are allowed
\usepackage{amsfonts} %Check we are allowed
\usepackage{booktabs} %Check we are allowed
\usepackage{graphicx} %Check we are allowed
\usepackage{tabularx}
\usepackage{float}
\usepackage{underscore} %Check we are allowed


%%% Colour commands - delete for final draft 
% - Editing Macros (Hierarchy of Importance) -
% 1. FACT CHECK (Critical): Red commands immediate attention for potential inaccuracies.
\newcommand{\fc}[1]{\textcolor{red}{[FactC: #1]}} 
% 2. CITE (High): Orange highlights missing evidence or attribution.
\newcommand{\ct}[1]{\textcolor{orange}{[CITE: #1]}} 
% 3. APPENDIX (Medium): Blue indicates structural moves (shifting text out of the main body).
\newcommand{\app}[1]{\textcolor{blue}{[APP: #1]}} 
% 4. SMALL REWORD (Low): Teal is for minor stylistic edits, mirroring your friend's setup.
\newcommand{\sr}[1]{\textcolor{teal}{[Sreword: #1]}}
% 5. FIGURE (Low): Purple is for figure-related notes (refer to figure, create figure, etc.).
\newcommand{\fig}[1]{\textcolor{purple}{[FIG: #1]}}
% 6. redraft (Low): fix immediadtely marker suggestions and corrections
\newcommand{\fiximm}[1]{\textcolor{brown}{[AI: #1]}}
% 7. LIMITATION (Low): Move to limiations 
\newcommand{\limit}[1]{\textcolor{cyan}{[LIM: #1]}}

%=======================================================================================
\begin{document}
\begin{preliminary}

\title{Physics-informed prediction and spatial attribution of forest top-height growth using repeated LiDAR surveys}
%TODO: DRAFT DIFFERENT VERSIONS OF THE TITLE AND PICK AGAIN AT THE END

\author{Shreeya Gorasia}

\date{\today}
%=======================================================================================

\abstract{TODO: FILL IN}

%TODO: writing: Comeback to this section when the entire report is done 
%=======================================================================================

\maketitle

\newenvironment{ethics}
   {\begin{frontenv}{Research Ethics Approval}{\LARGE}}
   {\end{frontenv}\newpage}

\begin{ethics}
This project was planned in accordance with the Informatics Research
Ethics policy. It did not involve any aspects that required approval
from the Informatics Research Ethics committee.

\standarddeclaration
\end{ethics}

%TODO: writing: Comeback to this section when the entire report is done 
%=======================================================================================

\begin{acknowledgements}
Any acknowledgements go here.
\end{acknowledgements}

%TODO: writing: Comeback to this section when the entire report is done 
%=======================================================================================

\tableofcontents
\end{preliminary}

%TODO: check: Comeback to this when entire report is done and check it and make it fits in one page
%=======================================================================================
\chapter{Introduction}  \label{ch:intro}

\paragraph*{Motivation and problem context}
Forest Digital Twins are a wider long-term goal for the forestry sector: digital models of forest systems that help monitor, forecast, and  management planning \cite{baskent2026Review_C1_SS, buonocore2022Proposal_C1_SS}. They aim to replace static, periodic checks with adaptive, data-driven planning by combining repeated field observations, simulation, and machine learning(ML), and artificial intelligence(AI) \cite{baskent2026Review_C1_SS, buonocore2022Proposal_C1_SS}. This is vital as managers face climate risks—like heavy storms, high winds, and severe summer heatwaves that spark wildfires and shift growth. At the same time, forestry is moving toward continuous cover management, which demands more flexible stand-level models \cite{ForestResearch2009Evidence_C1_SS}.

Sitka spruce (\textit{Picea sitchensis}) is Britain's primary commercial species, making up approximately 50\% of managed forests \cite{ForestResearch2026Sitka_C1_SS} and contributing roughly \pounds1.1 billion annually to Scotland's economy \cite{ScottishForestry2026Future_C1_SS, mason2011Sitka_C1_SS}. The study focuses on Aberfoyle in Trossachs, Scotland, a Sitka spruce-dominated region with 21 years of repeated LiDAR coverage. Featuring varied topography with exposed, windy ridges to sheltered valleys, providing variation that causes different growth trajectories in different areas of the forest.

% Height, LiDAR and shared growth relationship
Top height is a key measure of stand development, commonly used in forestry models to estimate downstream metrics like timber volume and site productivity \cite{ForestResearch2026Sitka_C1_SS}, and now can be accurately captured at scale using airborne LiDAR scans \cite{suarezminguez2008Practical_C1_SS}. For most tree species, including Sitka spruce, height development over time is modelled by the Chapman-Richards equation \cite{pommerening2015Methods_C2_CR}, a sigmoid growth curve, with its parameters fixed to represent a stand's maximum height potential and its speed of growth. A fixed parameter curve only captures the average population trend, but local trajectories can deviate from this baseline, due to factors like terrain, climate, soil, wind exposure, and active management all critically constrain or enhance a stand's actual development \cite{socha2023Higher_C1_SS, cameron2015Building_C1_SS}.  This raises two key questions: whether local departures from this shared curve are associated with identifiable environmental and management conditions, and whether independent attribution methods agree on which conditions matter.

Foresters need to know if a stand's assigned yield class still holds, or which local factors push growth off track. This dissertation tests whether letting these environment--growth relationships vary across space captures local growth patterns that fixed regional curves miss. Flexible machine learning models can also learn these patterns directly from data without fixed curve assumptions; we test deep neural networks (DNNs) and gradient-boosted trees (XGBoost) as a supplementary check on top-height predictions.

\paragraph*{Research gap and Contribution}
While repeated airborne laser scanning has been applied to model forest height growth and productivity \cite{janiec2023Development_C1_SS}, limited evidence exists on whether local environmental conditions are associated with departures from Sitka spruce's shared height-growth curve, and whether different attribution methods identify consistent conditions. Many current remote-sensing models do not account for spatial dependence between nearby-identical trees, leading to data leakage and inflated performance estimates \cite{Karasiak2022Spatial_C5_EVAL, ploton2020Spatial_C5_EVAL, lynch2025Digital_C4_NN}. Relying on black-box neural networks without explainable AI limits their use for stakeholders who require interpretable insights to support decision-making.

This dissertation covers one part of the Digital Twin idea: using attribution methods to explain local growth deviations, and checking if that explanation can help predict top height. First, it tests whether growth-environment links change by location, using a Geographically Neural Network-Weighted Regression (GNNWR) model. Second, it compares standard and physics-informed neural networks (PINNs) on prediction itself, to see which one wins and why.

\section{Objectives and Research Questions}

Answering the three research questions below required cleaning and combining repeated forest inventory and airborne LiDAR records with multi-source geospatial rasters (terrain, wind, climate, and soil) into a single unified modelling dataset (Chapter~\ref{ch:data}). Together, these questions form a practical, tiered workflow for forest management: Q1 flags stands that pull away from the standard regional growth curve and identifies the environmental drivers behind the split; Q2 tests whether a stand's assigned yield-class ceiling ($\Delta y_{\max}$) needs local updating; and Q3 evaluates whether adding these site features and physics constraints improves ongoing top-height prediction in the field (Table~\ref{tab:aims_objectives}).

\begin{table}[!htb]
\small
\setlength{\tabcolsep}{5pt}
\renewcommand{\arraystretch}{1.08}
\begin{tabularx}{\linewidth}{@{} >{\raggedright}p{0.32\linewidth} >{\raggedright\arraybackslash}X @{}}
\noalign{\hrule height 1pt}
\textbf{Research Question \& Target} & \textbf{Objectives \& Operational Purpose} \\
\noalign{\hrule height 0.6pt}
\textbf{Q1: Baseline Curve Attribution} \newline
\textit{Target:} Multi-year mean residual ($\bar{r}_i$) from the shared Chapman--Richards curve &
\textbf{(1.1)} Link persistent plot-level growth deviations to terrain, soil, wind, and management. \newline
\textbf{(1.2)} Test if drivers agree across Elastic Net, LMM, and XGBoost. \newline
\textbf{Purpose:} Flags off-track stands and tests attribution consistency across model types. \\
\hline
\textbf{Q2: Spatial Variation in Growth} \newline
\textit{Target:} Yield-class asymptote deviation ($\Delta y_{\max}^{(i)}$) &
\textbf{(2.1)} Test spatially varying models (GNNWR) under spatial CV to capture local patterns. \newline
\textbf{(2.2)} Separate genuine local environmental trends from noise and fitting artifacts. \newline
\textbf{Purpose:} Diagnose if a stand's standard yield class needs revision based on local site effects. \\
\hline
\textbf{Q3: Prediction \& Physics Guidance} \newline
\textit{Target:} Direct LiDAR top height ($H_{it}$) &
\textbf{(3.1)} Compare standard models against two physics-informed neural networks: PINN, which adjusts height only, and PINN-$k$, which adjusts both height and growth rate. \newline
\textbf{(3.2)} Test whether environment data or physics rules improve accuracy or make the fitted parameters easier to interpret, and weigh the trade-off between the two PINN versions. \newline
\textbf{Purpose:} Finds the best model for predicting future height, and whether PINN-$k$'s extra parameter is worth its small drop in accuracy. \\
\noalign{\hrule height 1pt}
\end{tabularx}
\caption{\small{Overview of research questions, objectives, and wider management purpose.}}
\label{tab:aims_objectives}
\end{table}


\paragraph*{Main Finding:}
Environmental and management conditions can explain some of the growth differences within Aberfoyle. Elastic Net, LMM, and XGBoost all point to similar features (Q1). GNNWR finds little that Elastic Net, LMM, and XGBoost miss, even though it lets the environment to growth link change across the forest. Most local patterns are just noise, or flaws in how the target was built (Q2). XGBoost is the most accurate model for predicting top height, outperforming the standard DNN and both PINN variants. While the two-parameter PINN loses a little accuracy, it ensures more biologically realistic results by never predicting impossible heights (Q3). In short, physics guidance and local spatial modelling do not improve accuracy, but they yield models that are safer and more interpretable.


%=======================================================================================
%=======================================================================================
\chapter{Related Work and Background} \label{ch:background}

\paragraph*{Shared growth models}
Top height tracks upper-canopy development and serves as a standard metric of forest growth and site productivity. Multi-temporal airborne LiDAR surveys now allow top height to be estimated consistently across broad landscapes, though uncertainty persists in young stands, open canopies, and edge boundaries \cite{ForestResearch2026Sitka_C1_SS, suarezminguez2008Practical_C1_SS}

The Chapman-Richards(CR) model describes height trajectories using a flexible sigmoid curve parametrised by an asymptotic ceiling ($y_{\max}$), an intrinsic growth rate ($k$), and a shape exponent ($p$). General Yield Class (GYC) groups stands by their expected timber production rate, with each class mapped to a standard height growth curve. While a single population curve offers a useful baseline, it cannot capture varying local site conditions and management histories \cite{chapman1961Statistical_C2_CR, pommerening2015Methods_C2_CR, mason2011Sitka_C1_SS}

\begin{figure}[!htbp]
    \centering
    \includegraphics[width=0.65\textwidth]
    {figures/fig_background/CR-curves-yeildclass.png}
    \caption{Chapman-Richards growth curves coloured by General Yield Class (GYC), work done by Forest Research}
    \label{fig:chapman-richards-illustrative}
\end{figure}

\paragraph*{Environmental controls}
Sitka spruce productivity responds strongly to elevation, wind exposure, soil moisture, and soil fertility. In upland northern Britain, General Yield Class falls by roughly $3.2\text{-}4.0\,\mathrm{m^3\,ha^{-1}\,yr^{-1}}$ per 100\,m rise in elevation, driven by lower temperatures and increased wind exposure \cite{worrell1990Productivity1_C1_SS}. Studies of Irish plantations similarly show peak productivity on moist, nutrient-rich sites, with sharp declines on drought-prone or waterlogged soils \cite{farrelly2011Sitka_C1_SS}.

Landscape-scale modelling integrates forest inventory data with spatial layers. Early UK research combined sub-compartment records with terrain and soil maps using principal component analysis and linear regression to predict yield classes \cite{bateman1998Using_C1_SS}. More recent studies pair LiDAR metrics with climate, topographic, and soil-type rasters. Because predictor importance varies across spatial resolutions and extents \cite{fricker2019More_C3_LID}, and coarse soil layers often correlate poorly with local soil conditions \cite{miller2024Assessing_C2_CR}, discrepancies between studies likely reflect differences in study extent, map resolution, and sampling design rather than a single universal list of environmental drivers.

\paragraph*{Global attribution}
Global models apply one environment-growth rule across the entire forest rather than letting it change by location. Mixed-effects models are common in forestry for separating broad environmental trends from stand-level differences \cite{harrison2018Brief_C5_EVAL, socha2016Assessment_C2_CR}. However, their random effects merely soak up leftover plot variance rather than explaining its ecological causes. Elastic Net handles correlated variables well by combining penalty terms ($L_1$ and $L_2$), but it still assumes relationships are strictly linear and additive \cite{zou2005Regularization_C4_NN}. Because each method relies on different mathematical assumptions, comparing these model families is a common way to test whether an association is robust, or an artefact of one model's own structure.

Boosted-tree models have been used to map forest productivity from LiDAR and environmental data, while XGBoost-SHAP has been applied to climatic and environmental drivers of forest condition \cite{puliti2023Tree_C2_CR, chen2016XGBoost_C4_NN}. When paired with tree models, SHAP has been used to identify non-linear thresholds and interactions learned by fitted tree models \cite{lundberg2017Unified_C7_XAI}. However, SHAP values reflect feature attribution within the model rather than direct biological causality, and attribution can be redistributed among correlated predictors \cite{aas2021Explaining_C7_XAI}. Limited work compares linear, mixed-effects, and tree-based attribution on the same persistent growth-departure target.

\paragraph*{Spatially varying attribution}
Spatial autocorrelation occurs when nearby measurements are similar to one another. By contrast, spatial non-stationarity means that the relationship between the environment and tree-growth changes from one geographic area to another. 

When a model leaves residual spatial autocorrelation, neighbouring areas are systematically over- or under-predicted. This may reflect missing spatial predictors, an oversimplified functional form, sampling effects, or genuine spatial non-stationarity \cite{gaspard2019Residual_C6_GWR}.

Geographically Weighted Regression (GWR) represents non-stationarity by fitting local regressions in which nearby observations receive greater weight. Its local coefficients can be mapped to show how estimated relationships vary across space \cite{brunsdon1996Geographically_C6_GWR}. However, local multicollinearity can produce unstable or strongly correlated coefficients, making these maps difficult to interpret \cite{wheeler2005Multicollinearity_C6_GWR}. Multiscale GWR (MGWR) relaxes GWR's single-scale assumption by allowing each predictor to operate over a different spatial bandwidth \cite{fotheringham2017Multiscale_C6_GWR}.

Geographically Neural Network-Weighted Regression (GNNWR) replaces predefined spatial kernels with a neural network that learns the relationship between distance and local weighting \cite{du2020Geographically_C6_GWR}. This addresses limitations of classical GWR, whose results depend on the selected bandwidth and kernel form and whose simple weighting functions may not represent complex spatial processes. In forestry, GNNWR improved prediction and reduced residual spatial autocorrelation when relating down wood to terrain and stand characteristics \cite{sun2022Local_C6_GWR}. However, evidence came from one forest reserve, one disturbance outcome, and five retained predictors, while the study gave limited attention to GNNWR's own interpretability and generalisability.


\paragraph*{Physics-guided prediction}
Physics-informed neural networks use governing equations as soft constraints during training \cite{raissi2019Physics_C4_NN}. Where exact physical laws are unavailable, process-guided models instead embed approximate biological knowledge. Forest applications have represented linked processes such as growth, competition, regeneration and mortality \cite{pichler2025Inferring_C4_NN}. However, rigid process assumptions can bias estimates of other processes, while sparse observations may leave alternative mechanisms difficult to distinguish. Even flexible hybrid models can learn biologically implausible responses in poorly sampled conditions \cite{pichler2025Inferring_C4_NN}.

Training also requires empirical and process losses to be balanced. Numerical stiffness can produce unequal gradients, causing one objective to dominate optimisation \cite{wang2021Understanding_C4_NN}. Chapman-Richards guidance therefore encourages an established growth shape but does not guarantee either biological validity or improved prediction. Although process-guided networks are widely studied, integrating Chapman-Richards equations into neural networks remains rare outside specialised variable-age growth models but it has not been tested under long-term LiDAR time series in Aberfoyle.

\paragraph*{Residual spatial evaluation}
Residual Moran's $I$ is used in ecology to spot spatial patterns that average error scores can miss. It gives a simple test of whether nearby plots are consistently over- or under-predicted, indicating if a model has missed key spatial structure or yields misleading conclusions \cite{moran1950Notes_c5_eval,kim2021Predicting_C5_EVAL}. In forestry, it helps compare global, GWR, and neural spatial models, where local neural models produced both lower errors and less residual clustering \cite{sun2022Local_C4_NN}. However, Moran's $I$ relies on the chosen distance settings and compresses the whole map into a single number, which can hide distinct local patches. A lower score shows less error clustering under those settings, but it cannot explain why or prove that the model found a true biological relationship.

\section{Research Gap and Contributions}
Existing studies show that the environment strongly shapes forest growth, but they typically treat growth curves, driver attribution, spatial variation, and prediction as separate problems. No previous study directly compares linear, mixed-effects, non-linear, and spatial attribution models using multi-temporal LiDAR trajectories on the same dataset with spatially independent validation. As a result, it remains unclear whether identified drivers agree across model types, whether their effects change across the landscape, and whether adding this environmental information actually improves top-height predictions.

This work addresses these gaps through three contributions:
\begin{enumerate}
\setlength\itemsep{0pt}
\item \textbf{Global attribution (Q1):}
Compares whether environmental associations with persistent curve departure remain consistent across Elastic Net, mixed-effects and XGBoost models.
\item \textbf{Spatial attribution (Q2):}
Tests whether associations with fitted asymptotic departure vary geographically using GNNWR and spatially blocked evaluation.
\item \textbf{Prediction and process guidance (Q3):}
Compares standard and neural top-height models, then tests whether environmental inputs or Chapman--Richards guidance improve accuracy or parameter interpretation.
\end{enumerate}
%=======================================================================================
%=======================================================================================

%=======================================================================================
%=======================================================================================
\chapter{Data}
This chapter describes the dataset used for model training and evaluation, outlining key feature engineering steps and data limitations. Table~\ref{tab:data_summary} summarises all predictor and target variables alongside their sources and temporal coverage. Feature selection rationale, data partitioning, and specific modelling choices are detailed in Chapter~\ref{ch:methodology}.

\section{Study area and cohort coverage}
The dataset covers approximately $16\text{~km} \times 12\text{~km}$ of Aberfoyle and contains 795,645 plot--survey records from 230,359 unique plot locations across six survey years (2002--2023). The records contain airborne laser scanning (ALS)-derived variables, including canopy cover and height percentiles, alongside stand and management fields such as planting year and known felling history (Table~\ref{tab:data-fields}). Figure~\ref{fig:data-study-area} shows the study-area topography, the spatial coverage of both cohorts, and an example of standard and boundary-clipped plots.

\begin{figure}[!htbp]
\includegraphics[
    width=1.05\textwidth,
    height=0.48\textheight,
    keepaspectratio
]{figures/fig_data/aberfoyle_area.png}
\caption{\small{Left:Aberfoyle study area. Middle: cohort coverage, Right: plot-level detail}}
\label{fig:data-study-area}
\end{figure}

Because LiDAR coverage increase by $37\text{ km}^2$ between 2006 and 2008, two nested cohorts were defined (Figure~\ref{fig:tab:data-cleaning-funnel}): a broad-coverage \emph{4-survey} cohort (2008, 2012, 2021, 2023; 71,766 plots across 296 compartments) and a longer-trajectory \emph{6-survey} cohort (2002--2023; 13,897 plots across 48 compartments).


To ensure data quality, sequential filters were applied to both cohorts (Table~\ref{tab:data-cleaning-funnel}):
\begin{itemize}
    \item \textbf{Complete presence:} Plots must be present across all survey years in the cohort.
    \item \textbf{Species purity:} Pure Sitka spruce stands only.
    \item \textbf{Plausible records:} Planting year between 1850 and the survey year, stand age $0\text{--}200$ years, and LiDAR top height $0\text{--}60\text{ m}$ (ecological maximum \cite{mason2011Sitka_C1_SS}).
    \item \textbf{Maturity cutoff:} Stand age $\ge 30$ years by 2023 to exclude unreliable young-stand height measurements \cite{suarezminguez2008Practical_C1_SS}.
\end{itemize}

\begin{table}[!htbp]
\centering
\scriptsize
\begin{tabular}{lrrrr}
\toprule
& \multicolumn{2}{c}{4-survey} & \multicolumn{2}{c}{6-survey} \\
\cmidrule(lr){2-3}\cmidrule(lr){4-5}
Stage & Plots & Removed & Plots & Removed \\
\midrule
0. Raw dataset
& 230,359 & -- & 230,359 & -- \\
1. Cohort years present
& 230,345 & 0.01\% & 230,359 & 0.00\% \\
2. Complete survey coverage
& 112,782 & 51.04\% & 18,781 & 91.85\% \\
3. Sitka spruce only
& 72,395 & 35.81\% & 13,897 & 26.01\% \\
4. Valid planting year
& 71,766 & 0.87\% & 13,897 & 0.00\% \\
5. Plausible age (0--200 years)
& 71,766 & 0.00\% & 13,897 & 0.00\% \\
6. Plausible height (0--60 m)
& 71,766 & 0.00\% & 13,897 & 0.00\% \\
7. Mature stands (age $\geq 30$ in 2023)
& \textbf{58,112} & 19.03\%
& \textbf{13,769} & 0.92\% \\
\bottomrule
\end{tabular}
\caption{\small{Plots left after each filter. Percentages show the share removed from the previous row, rather than from the raw dataset. Figure~\ref{fig:data-study-area} shows the final spatial coverage.}}
\label{tab:data-cleaning-funnel}
\end{table}

%------------------------------------------------
\section{Response variable and derived attribution targets}
All targets in this dissertation derive from the plot-level 95th percentile ALS return height (\texttt{elev\_percentile\_95th}), forming three distinct response variables (detailed in Chapter~\ref{ch:methodology}): the multi-year mean residual from a shared growth curve (Q1, $\bar{r}_i$), the gap between a plot's fitted ALS height ceiling and its assigned yield-class ceiling (Q2, $\Delta y_{\max}^{(i)}$), and direct LiDAR top height (Q3, $H_{it}$). Forest Research records supply planting year, management history, and yield class (Table~\ref{tab:data-fields}), with any variables calculated directly from measured height removed from model inputs to prevent target data leakage. A code-level audit confirmed that the pipeline contains no target circularity: environmental features (terrain, climate, wind, and soil) were extracted independently from external spatial layers and joined only at the modelling stage, never entering the target calculation itself. However, because historical yield-class assignments come directly from Forestry Commission records, the original field assessments that generated them could not be independently verified.


\begin{table}[!htbp]
\scriptsize
\setlength{\tabcolsep}{3pt}
\renewcommand{\arraystretch}{0.95}
\begin{tabular}{
    >{\raggedright\arraybackslash}p{2.6cm}
    >{\raggedright\arraybackslash}p{7.0cm}
    >{\raggedright\arraybackslash}p{3.2cm}
}
\hline
Field & Meaning or derivation & Role in this study \\
\hline
\texttt{elev\_\allowbreak percentile\_\allowbreak 95th} & 95th percentile of ALS return heights above ground, used without correction & Response variable \\
\texttt{Top\_Height95} & $=\texttt{elev\_percentile\_95th}\times1.1$; Forest Research's own corrected height estimate & Retained for reporting but not modelled \\
\texttt{plyr}, \texttt{Age} & Planting year; \texttt{Age} = survey year $-$ \texttt{plyr} & Input feature; used for data cleaning \\
\texttt{yldc} & An existing Forest Research yield-class field, not calculated from this project's LiDAR-derived height response & Excluded from prediction inputs; used to construct Q2's own attribution target \\
\texttt{whcl} & Ordinal Forest Research windthrow hazard class, 0--6 & Management descriptor \\
\texttt{CanopyCover} & Canopy-cover fraction from the same ALS survey and processing step as the height response & Input feature \\
\texttt{Thin}, \texttt{last\_thinn}, \texttt{next\_thin\_} & Thinning history; \texttt{Thin} is binary (1 if thinned, 0 if never thinned); \texttt{last\_thinn}/\texttt{next\_thin\_} record last and next thinning years & Input feature \\
\texttt{Vol95}, \texttt{area}, \texttt{g1}, \texttt{g3} & Forest Research volume estimate, calculated from height, plot area, and canopy cover & Used for disturbance screening \\
\texttt{GYCspec95} & General Yield Class calculated directly from a plot's own measured height (\texttt{Top\_Height95}) & Excluded (target leakage) \\
\texttt{p1}--\texttt{p5} & Forest Research growth curve parameters for yield class & Excluded from inputs; used with \texttt{yldc} to build Q2 target curve \\
\hline
\end{tabular}
\caption{\small{Target, stand, and management fields provided by Forest Research. \texttt{GYCspec95} and \texttt{p1}--\texttt{p5} are excluded from model inputs to prevent target data leakage. In Q2, \texttt{p4} and \texttt{p5} are held fixed only to define the baseline yield-class asymptote target.}}
\label{tab:data-fields}
\end{table}

%------------------------------------------------
\section{Derived terrain, wind, climate and soil variables}
\label{sec:data-environmental}

Each plot was matched with external terrain, wind, climate, and soil variables (Table~\ref{tab:data-env-sources}). Appendix~\ref{tab:data-env-variables_detail_app} defines each variable.

\begin{table}[!htbp]
\centering
\scriptsize
\setlength{\tabcolsep}{4pt}
\renewcommand{\arraystretch}{1.05}
\begin{tabular}{ >{\raggedright\arraybackslash}p{2.0cm} >{\raggedright\arraybackslash}p{4.2cm} >{\raggedright\arraybackslash}p{2.3cm} >{\raggedright\arraybackslash}p{4.0cm}}
\hline
Data family & Source and resolution & Time basis & Variables used \\
\hline
Terrain & OS Terrain 50, 50\,m \cite{ordnancesurveyTerrain50_C3_LID} & Static & Elevation, slope, aspect, northness and eastness \\

Soil & CEH Natural Capital, 50\,m \cite{feeney2026CEHNaturalCapital_C3_LID}; SoilGrids v2.0, 250\,m \cite{poggio2021SoilGrids_C3_LID} & Static & Pedotope, drainage, texture and soil pH \\

Wind exposure & Global Wind Atlas, 250\,m \cite{davis2023GlobalWindAtlas_C3_LID} & 2006--2015 climate & Long-term wind speed and exposure \\

Observed wind & MIDAS Open weather stations \cite{metoffice2019MIDASOpen_C3_LID} & Survey intervals & Mean and extreme wind conditions \\

Long-term climate & CHELSA v2.1, 1\,km \cite{karger2017CHELSA_C3_LID} & 1981--2010 climate & Growing degree days and precipitation \\

Matched climate & HadUK-Grid, ~1\,km \cite{hollis2019HadUKGrid_C3_LID} & Cohort years & Temperature and frost \\

Spatial context & OS Open Roads and Open Rivers, vector data
\cite{ordnancesurveyOpenRoads_C3_LID,ordnancesurveyOpenRivers_C3_LID} & Static & Distance to roads, rivers and forest edges \\
\hline
\end{tabular}

\caption{\small{External data sources used in the analysis. Values were matched to plot centres. Full source, resolution, coverage and extraction details are given in Appendix~\ref{app:data-sources}.}}
\label{tab:data-env-sources}
\end{table}

Environmental variables vary in their time resolution. Terrain, soil, climate normals (CHELSA, Global Wind Atlas), and TOPEX were kept static, assigning each plot a single value across all survey years. TOPEX measures topographic shelter, where higher scores mean greater shelter and lower scores mean higher exposure. By contrast, MIDAS wind metrics were calculated specifically for each interval between survey dates using weather-station data. HadUK-Grid temperature and frost values were averaged over the survey years within each cohort, while road, river, and edge distances were treated as static baselines. Limited missing values in environmental features were handled by removing incomplete rows for specific model sets or using missingness indicators for thinning history (Chapter~\ref{ch:methodology}).

Several limitations should be noted. Requiring plots to appear across every survey year ensured clean tracking over time, but reduced landscape coverage, particularly in the 6-survey cohort where high-elevation areas are underrepresented. Predictor layers also vary in detail: terrain changes across individual plots, whereas regional climate maps share the same values across wide areas and do not update per survey date. Additionally, canopy cover shares the same LiDAR flight acquisitions as top height, which could introduce shared survey conditions. These limits are especially relevant when evaluating local plot-level attribution.


%=======================================================================================
%=======================================================================================
\chapter{Methodology and Experimental Framework} \label{ch:methodology}
\section{Experimental framework}
Q1 links a plot's deviation from the baseline growth curve to environmental and management factors. Q2 tests whether place-specific relationships explain patterns that global models miss. Q3 tests whether these signals improve top-height prediction. 
\begin{table}[!ht]
\centering
\footnotesize
\renewcommand{\arraystretch}{1.2}
\begin{tabularx}{\linewidth}{@{} l >{\raggedright\arraybackslash}X >{\raggedright\arraybackslash}X >{\raggedright\arraybackslash}X @{}}
\toprule
Q & Unit and target & Models or comparison & Main evaluation \\
\midrule
Q1 & Plot; how far a stand sits from the shared growth curve, averaged over its survey history ($\bar r_i$) -- positive means taller than expected, negative means shorter & Elastic Net, XGBoost and a linear
mixed-effects model & Spatial block CV (Elastic Net, XGBoost); LMM coefficients and variance decomposition; residual Moran's $I$ \\
Q2 & Plot; gap between a plot's own fitted height ceiling and the ceiling its yield class predicts ($\Delta y_{\max,i}$) -- positive means a higher ceiling than expected & Elastic Net, XGBoost
and GNNWR & Five-fold spatial block CV; attribution, Moran's $I$ and target audit \\
Q3 & Plot--survey row; top height itself ($H_{it}$) & CR, linear regression, random forest, XGBoost, DNN, PINN and PINN-$k$
& Five-fold spatial block CV; pooled accuracy metrics; sensitivity checks on training
settings, physics-loss weight and environmental feature-set size; permutation-importance
and age/height error diagnostics \\
\bottomrule
\end{tabularx}
\caption{\small{Units, targets and evaluation used for each question.}}
\label{tab:rq-framework}
\end{table}


\paragraph*{Data partitions and leakage control}
The 4-survey cohort served as our primary dataset, while the smaller 6-survey cohort provided a longer-term sensitivity check. For Q1 and Q2, each plot was reduced to a single summary value; for Q3, we used every individual plot-survey measurement over time. To prevent data leakage, we split the data by forest compartment rather than by individual plot \cite{roberts2017Cross_C5_EVAL}. No compartment appeared in both training and testing sets. We also removed any training plots within 60\,m of validation or test stands to cut down on spatial spillover. We grouped the compartments into five equal folds. In each round, three folds trained the model, one tuned its settings, and one evaluated final accuracy. Across the five rounds, every compartment was tested exactly once, and all test predictions were pooled to compute the overall scores \cite{ploton2020Spatial_C5_EVAL}.

\begin{figure}[H]
\centering
\includegraphics[width=0.85\textwidth]{figures/fig_methodology/fig_method_map_splits.png}
\caption{\textbf{Data-splitting strategies for model evaluation.} 
(A)~Random plot-level split: shown for comparison only, as spatial proximity between adjacent training and test plots caused severe data leakage and inflated accuracy. 
(B)~Compartment-level spatial-block split with a 60\,m buffer: the primary strategy used throughout this dissertation to ensure complete spatial independence between training, validation, and test sets.}
\label{fig:splits-combined}
\end{figure}

Continuous variables were normalized using training-set statistics and applied to validation and test folds. Stand age was scaled separately so PINN outputs could be directly converted back to metres per year. Soil classes were one-hot encoded, while administrative IDs (compartment and block) were excluded from all feature sets. Additionally, two spatial neighbour-height variables were dropped because they were computed from adjacent plot heights prior to splitting and would have leaked test-set information.

%------------------------------------------------
\paragraph*{Environmental feature sets}
Four feature sets were constructed to test whether adding environmental data improves model performance. Because feature relevance depends on the specific prediction task, sets were ranked and selected separately for Q1, Q2, and Q3. Variables with extensive missing values (e.g., $>20\%$) and exact duplicates were removed first. For highly correlated pairs (Spearman $|\rho| \geq 0.95$), we retained only the directly measured variable. The remaining candidates were ranked using their average rank across three metrics evaluated on a single spatial split: Spearman $|\rho|$, XGBoost permutation importance, and XGBoost drop-column ablation. To prevent multicollinearity, we iteratively removed variables with a variance inflation factor (VIF) greater than 5 until all remaining features met the threshold \cite{obrien2007Caution_C5_EVAL}. Baseline stand variables were retained unconditionally but still included in VIF calculations. This screening removed 430 four-survey plots, primarily due to gaps in SoilGrids coverage; the six-survey cohort had complete environmental coverage. Full missingness summaries and exact feature lists are provided in Appendix~\ref{app:feature-sets}.

\begin{table}[H]
\centering
\scriptsize
\renewcommand{\arraystretch}{1.2}
\begin{tabularx}{\linewidth}{@{}p{1.4cm}X@{}}
\toprule
Set & Construction \\
\midrule
Set~1 & Stand and management baseline only. \\
Set~2 & Set~1 plus the top 10 ranked environmental candidates (candidates failing VIF screening were replaced by the next-ranked variable). \\
Set~3 & Set~1 plus the top 50\% of terrain and wind features, followed by VIF screening. \\
Set~4 & Set~3 candidates plus the top 50\% of climate, soil, and spatial-context features, followed by VIF screening across the pooled set. \\
\bottomrule
\end{tabularx}
\caption{\textbf{Construction of environmental feature sets.} Feature ranking was performed independently for Q1, Q2, and Q3; full variable listings are given in Appendix~\ref{app:feature-sets}.}
\label{tab:feature-sets-construction}
\end{table}

Because VIF screening was rerun whenever the candidate pool expanded, the final sets are not strictly nested. For example, elevation was selected in Sets~2 and~3 but dropped from Set~4 due to collinearity with broader terrain and climate predictors. This reflects statistical redundancy within the larger pool rather than an absence of ecological influence.
%------
\section{Q1: Explaining growth deviations}
Q1 evaluates which forest stands consistently outperform or lag behind the regional Chapman--Richards growth curve once age is accounted for, and identifies the environmental drivers behind these differences. The analysis proceeds in two steps: first, repeated survey heights are condensed into a single mean residual per plot; second, Elastic Net, XGBoost, and a linear mixed-effects model are trained to predict this deviation. We use the 4-survey cohort because the 6-survey cohort contains only 47 compartments, which is insufficient for robust spatial-block cross-validation.

\paragraph*{What Q1's target means in plain terms}
Let $H_{it}$ and $a_{it}$ denote the observed top height and stand age for plot $i$ at survey $t$, where $T_i$ is the total number of surveys for plot $i$. Using the regional baseline height $H(a_{it})$ from Equation~\eqref{eq:cr}, the survey-level residual $r_{it}$ and plot-level mean departure $\bar{r}_i$ are defined as:
\begin{align}
r_{it} &= H_{it} - H(a_{it}), \label{eq:cr-gap} &
\bar{r}_i &= \frac{1}{T_i} \sum_{t=1}^{T_i} r_{it}. \label{eq:mean-cr-residual}
\end{align}
A positive $\bar{r}_i$ indicates that a plot tracks above the regional baseline on average, while a negative value indicates under-performance. Because $\bar{r}_i$ reflects average vertical elevation relative to the curve rather than instantaneous growth rate, a plot's residual sign may vary across individual survey years. Stand age is excluded from the feature set since the baseline curve already removes the age trend.

\paragraph*{Models}
Both Elastic Net and XGBoost are evaluated using 5-fold spatial-block cross-validation. Elastic Net tunes its regularization penalty and $\ell_1$-ratio via inner 5-fold cross-validation (\texttt{ElasticNetCV}). XGBoost uses fixed hyperparameters (max depth 4, learning rate 0.04, up to 500 trees with early stopping), adopted after initial validation checks rather than exhaustive re-tuning on this specific target (the direct top-height XGBoost in Q3 uses its own targeted grid search). Model performance is reported using out-of-fold pooled metrics ($R^2$, RMSE, and MAE). 

To account for plots grouped within compartments, a linear mixed-effects model is fit on each fold's training partition:
\begin{equation}
\bar{r}_i = \beta_0 + \boldsymbol{\beta}^{\top}\mathbf{x}_i + u_{c[i]} + \varepsilon_i, \qquad u_c \sim \mathcal{N}(0, \sigma_u^2).
\label{eq:mixed-effects}
\end{equation}
Here, $\boldsymbol{\beta}^{\top}\mathbf{x}_i$ is the fixed-effects term. It describes broad relationships across all plots, such as slope or TOPEX. $u_{c[i]}$ is the random intercept for plot $i$'s compartment. Each training compartment gets its own fitted intercept, drawn from a shared distribution with variance $\sigma_u^2$. This term captures unexplained grouping at compartment level. It does not identify the specific cause of that grouping. This model is used only for fitted coefficients and the variance-explained comparison below. It is not used to make held-out predictions for unseen compartments.

Parameters ($\boldsymbol{\beta}$, $\sigma_u^2$) are fitted with restricted maximum likelihood (REML). This avoids the small variance underestimate that plain maximum likelihood (ML) can give. When models with different fixed effects are compared, ML is used instead because REML likelihoods cannot be compared across different fixed-effect structures. This compares a null model (compartment random intercept only) against the full model (fixed effects added); variance explained = share of the null's compartment-level variance that shrinks once fixed effects are included. The model is refit on each fold's training data, so coefficients are reported as mean $\pm$ SD across folds. This shows fold-to-fold stability. It is different from the pooled Elastic Net and XGBoost metrics above, which measure held-out accuracy.

\paragraph*{Sensitivity checks}
Two concerns motivate four checks. First, canopy cover and top height come from the same LiDAR flights. A strong canopy-cover effect could therefore reflect measurement overlap rather than real stand structure. Three steps test this: (1) refit Set~4 without canopy cover and check the change in $R^2$; (2) check the direct correlation between canopy cover and each target; (3) recompute SHAP values without canopy cover to see if other predictors take over its role.

Second, terrain coefficients for slope and TOPEX sometimes differ between Elastic Net and the mixed-effects model. This could mean local compartment effects that global models miss. A fourth check measures how much slope and TOPEX variance sits between compartments versus within them (the intraclass correlation coefficient, ICC).

Elastic Net coefficient signs, mixed-model fixed-effect signs, and XGBoost mean $|$SHAP$|$ rankings are then compared. This shows which variables matter, and in which direction when sign is available. The raw sizes are not directly comparable because the models differ, and mean $|$SHAP$|$ has no sign. Agreement across all three makes an association less likely to be a quirk of one model, but it does not show causality. Residual Moran's $I$ checks for remaining spatial clustering (Section~\ref{sec:eval-robustness}).
%------------------------------------------------
\section{Q2: Spatial attribution}
Q2 builds on Q1, instead of asking only which variables matter overall, it asks whether their effects change by place. Global models assume one environment--growth relationship across the study area. Q2 tests whether allowing this relationship to vary by location improves explanation and reduces residual spatial clustering \cite{du2020Geographically_C6_GWR}. It gives each plot a prediction that depends on its location, based on the idea that nearby plots behave more similarly than distant ones. Its coefficients can therefore vary across the map, a neural network learns these spatial weights directly from geographic coordinates, without a fixed neighbourhood size chosen in advance. GNNWR still assumes that each local relationship is a straight line. If the real biological response is curved, a changing local coefficient may simply show that the local line fits badly there. It may not reflect a true difference in growth drivers by place. For that reason, both model accuracy and local coefficients are considered together.


\paragraph*{What Q2's target means in plain terms} Q2 asks whether a plot grows toward a higher or lower final height than its official yield class predicts. $p_1$--$p_5$ are Forest Research's published numbers that convert a yield-class rating into a full growth curve for a species and region. This project did not derive them, just reused them. $p_1$--$p_3$ set the curve's asymptote (ceiling); $p_4$--$p_5$ set its shape (rate and curvature) -- two parts of one sigmoid. $p_4$ and $p_5$ are fixed at their published values for the matching Forest Research curve; only the height ceiling is allowed to differ from plot to plot. Earlier work has represented environmental variation through the Chapman--Richards asymptote in the same way \cite{scolforo2013Dominant_C2_CR}. This supports it as a simple and workable choice. For plot $i$ at survey $t$:
\begin{align}
g_{it} &= \left(1-e^{-p_{4,i}a_{it}}\right)^{p_{5,i}}, \label{eq:git-curve} \\
\hat y_{\max,i} &= \frac{\sum_t H_{it}g_{it}}{\sum_t g_{it}^2}, \label{eq:plot-ymax-fit} \\
y_{\max,i}^{\mathrm{yldc}} &= \mathrm{yldc}_i p_{2,i}+p_{1,i}+2p_{3,i}, \label{eq:yldc-asymptote} \\
\Delta y_{\max,i} &= \hat y_{\max,i}-y_{\max,i}^{\mathrm{yldc}}. \label{eq:local-ymax-diff}
\end{align}
With $p_4$ and $p_5$ fixed, Equation~\eqref{eq:plot-ymax-fit} gives a least-squares estimate for $\hat y_{\max,i}$. This avoids trying to estimate three correlated curve parameters from only 4--6 observations per plot. $\hat y_{\max,i}$ is the ceiling fitted from plot $i$'s own observed trajectory, $y_{\max,i}^{\mathrm{yldc}}$ (Equation~\eqref{eq:yldc-asymptote}) is the ceiling implied by its assigned yield class, using the same fixed rate and shape. Their difference, $\Delta y_{\max,i}$ (Equation~\eqref{eq:local-ymax-diff}), is positive when the fitted ceiling is higher than expected and negative when it is lower. This target is estimated, not directly observed.

Following the Q1 setup, Elastic Net and XGBoost are trained on $\Delta y_{\max}$ using Sets~2--4, with the same tuning approach as Q1. Elastic Net tests a global linear relationship. XGBoost tests a global non-linear relationship. GNNWR tests a spatially varying but locally linear relationship. Spatial variation is modelled using GNNWR via its Python library \cite{du2020Geographically_C6_GWR, du2022PyGNNWR_C6_GWR}:
\begin{equation}
\widehat{\Delta y}_{\max,i} = \beta_0(s_i) + \sum_j \beta_j(s_i) x_{ij}.
\label{eq:gnnwr}
\end{equation}
For each location $s_i$, GNNWR uses a neural network to learn spatial weights from distances to training plots. These weights determine the local coefficients $\beta_j(s_i)$ around a global linear baseline. Each training fold uses only that fold's own training partition. No hyperparameter search was performed for GNNWR. Every run uses the library's default architecture and training settings.


\subsection*{Computational adjustments and feature evaluation}
Full reference sets were used for the headline GNNWR runs. Early preliminary tests had used a smaller capped reference set because of GPU memory limits. That cap is kept only as a later sensitivity check. Memory-intensive hat-matrix diagnostics were disabled. As a result, GNNWR AIC is approximated using raw feature counts rather than the true GWR-corrected effective degrees of freedom, and $F$-tests are not reported. Predictive metrics ($R^2$, RMSE) are unaffected. Elastic Net and XGBoost were unaffected. Feature importance is measured with XGBoost gain, test-fold SHAP values \cite{lundberg2017Unified_C7_XAI}, and GNNWR's local coefficient distributions and maps. Each feature group's contribution to $R^2$ is measured by permutation importance, by shuffling that group's columns without refitting. Its contribution to residual clustering is measured by removing the group from Set~4 and refitting, then comparing Moran's $I$ before and after on the four-survey cohort. A separate check isolates canopy cover. Each of the three models is refit with canopy cover as the only predictor, and again with canopy cover removed, to test how much of the score depends on this one variable.

\subsection*{Target and outlier audit}
Plots affected by harvesting or clear measurement errors are screened using two rules across consecutive surveys. First, clearfelling is flagged when top height drops $\ge 25\%$, canopy cover drops $\ge 70\%$, and volume drops $\ge 70\%$. Second, measurement error is flagged when top height drops $\ge 25\%$, but canopy cover drops $\le 10\%$ and volume does not fall. Plots with large height drops that match neither rule are kept. Extreme fitted height ceilings ($\hat{y}_{\max,i}$) are flagged using Tukey's $1.5 \times \text{IQR}$ rule \cite{tukey1977Exploratory_C5_EVAL}:
\begin{equation}
\left[Q_1 - 1.5 \times \text{IQR},\; Q_3 + 1.5 \times \text{IQR}\right].
\label{eq:tukey-fences}
\end{equation}
These bounds are calculated directly from the cohort data rather than using a fixed external cutoff. Flagged plots are kept in the main models, but evaluated in a separate check with those plots removed from test scores. Top 1\% and 5\% largest-residual plots are compared for overlap across Elastic Net, XGBoost, and GNNWR, to see whether the same plots are flagged by every method. Their spatial spread, distance to stand edges, wind shelter, and local coefficients are compared against the rest of the plots, to check whether extreme errors are linked to fitting artifacts, disturbance, or environmental conditions -- a link, not a proven cause. Which plots are flagged, and what that means, is reported in Results.

%------------------------------------------------
\section{Q3: Prediction and physics guidance}
Q3 compares models for predicting top height: a baseline growth curve, standard tabular models, and neural networks, including physics-informed ones. It then tests whether environmental variables or Chapman--Richards(CR) guidance improve accuracy or the interpretability of the fitted parameters. Using five spatial folds on Set~3, the other feature sets test whether adding environmental data helps.

\paragraph*{CR baseline}
\begin{equation}
H(a) = y_{\max}\left(1 - e^{-ka}\right)^p,
\label{eq:cr}
\end{equation}
where $H(a)$ is height at age $a$, $y_{\max}$ is the height ceiling, $k$ is the growth rate, and $p$ controls curve shape. All three parameters are fit freely by non-linear least squares on each training fold and then frozen, keeping test folds fully independent.

\paragraph*{Conventional models}
We test three baselines: linear regression, random forest, and XGBoost [1]. For Q3, XGBoost was tuned across 27 parameter settings over five folds using validation $R^2$ (Appendix~\ref{app:hyperparameters}).

\paragraph*{DNN and physics-informed models}
All neural networks: three-layer trunk, 128 units per layer, Leaky-ReLU, Adam optimiser (learning rate $10^{-4}$, weight decay $10^{-5}$, batch size 256, early stopping) -- defaults, checked against alternatives (Sensitivity checks and diagnostics, below). Standard DNN: minimises mean squared error ($\mathcal{L}_{\mathrm{data}} = \frac{1}{N}\sum (\hat H_i-H_i)^2$). Conditioned DNN: environmental variables feed the input directly. PINN, PINN-$k$: environmental variables feed small sub-networks instead, which adjust each plot's own growth-curve parameters. Plain PINN adjusts only the height ceiling ($y_{\max}$); growth rate ($k$) and shape ($p$) stay one shared value for every plot. PINN-$k$ adjusts both $y_{\max}$ and $k$; $p$ stays shared for both variants. These plot-level parameters adjust the predictions and loss terms together, balancing fit to the data with the baseline growth curve.

Equations~\eqref{eq:cr-derivative}--\eqref{eq:total-loss}: the physics-informed loss.
\begin{align}
H'(a) &= y_{\max}pk e^{-ka}\left(1-e^{-ka}\right)^{p-1}, \label{eq:cr-derivative} \\
\mathcal{L}_{\mathrm{phys}} &= \frac{1}{N}\sum_{i=1}^{N}\left(\frac{\partial\hat H_i}{\partial a_i}-H'(a_i)\right)^2, \label{eq:physics-loss} \\
\mathcal{L}_{\mathrm{traj}} &= \frac{1}{M}\sum_{j=1}^{M}\left(\frac{\hat H_{j,\mathrm{later}}-\hat H_{j,\mathrm{earlier}}}{\Delta a_j}-H'(\bar a_j)\right)^2, \label{eq:trajectory-loss} \\
\mathcal{L} &= \mathcal{L}_{\mathrm{data}} + \lambda_{\mathrm{phys}}\mathcal{L}_{\mathrm{phys}} + \lambda_{\mathrm{traj}}\mathcal{L}_{\mathrm{traj}} + 10^{-5}\sum_{\theta}|\theta|. \label{eq:total-loss}
\end{align}
Equation~\eqref{eq:cr-derivative}: theoretical growth rate, shown with the shared global parameters; conditioned models substitute each plot's own adjusted parameters instead, $y_{\max}^{(i)}$ for both PINN variants, plus $k^{(i)}$ for PINN-$k$. Equation~\eqref{eq:physics-loss}: gap between this rate and the network's own derivative ($\partial\hat H_i/\partial a_i$). Equation~\eqref{eq:trajectory-loss}: gap between consecutive-survey growth and this rate, at midpoint age $\bar a_j$. Equation~\eqref{eq:total-loss}: data loss plus weighted physics and trajectory terms, plus an L1 penalty. Defaults: $\lambda_{\mathrm{phys}}=\lambda_{\mathrm{traj}}=1$; 0 turns the physics terms off.

Sub-networks adjust plot parameters:
\begin{equation}
y_{\max}^{(i)} = y_{\max}+f_y(\mathbf{x}_i), \quad k^{(i)} = k\exp\!\left(f_k(\mathbf{x}_i)\right).
\label{eq:param-adjust}
\end{equation}
Plain PINN: $y_{\max}^{(i)}$ line only, $k^{(i)}=k$ fixed. PINN-$k$: both lines. Each sub-network in use ($f_y$ alone for PINN; $f_y$, $f_k$ for PINN-$k$): one hidden layer, 16 units, turns terrain features into one number. Trunk and sub-networks train together: one Adam optimiser, one learning rate, one weight decay, no separate tuning step per part. Sub-network outputs feed the prediction directly (Equation~\eqref{eq:pinn-prediction}), so ordinary training already shapes them; $\lambda_{\mathrm{phys}}$ adds a further pull toward matching the theoretical growth rate.

The final prediction: the CR curve at the plot's adjusted parameters, plus a residual from the trunk:
\begin{equation}
\hat H_i = y_{\max}^{(i)}\left(1-e^{-k^{(i)}a_i}\right)^{p} + \text{trunk residual}.
\label{eq:pinn-prediction}
\end{equation}
$y_{\max}^{(i)}$ shapes the prediction directly for both variants; $k^{(i)}$ does too, for PINN-$k$. Implemented in \texttt{pinn\_env\_terrain\_fix.py} / \texttt{pinn\_env\_terrain\_k\_fix.py}.

\subsection*{Sensitivity checks and diagnostics}
Beyond five-fold spatial evaluation, six further checks test the PINN and PINN-$k$ results. Each is a single run per fold or condition, not a repeated-seed study:
\begin{itemize}
\setlength\itemsep{0pt}
\item Training settings and physics weight: Alternative learning-rate and weight-decay settings, an alternative batch size (512 vs.\ 256), and eight $\lambda_{\mathrm{phys}}=\lambda_{\mathrm{traj}}$ settings (0-3.0) were each tested for DNN, PINN, and PINN-$k$ over five folds, against the defaults used for every other number in this chapter.
\item Environmental feature-set size: PINN and PINN-$k$ were retrained with no environment, and with Sets~2, 3, and 4, over five folds each. This tests whether environmental data helps at all, and whether the exact list matters.
\item Terrain permutation importance: Each terrain or wind feature was shuffled in turn, with ten repeats averaged. This measures how much personalised $y_{\max}$ and $k$ change.
\item Age- and height-banded error: MAE was compared by age band and height band for DNN, PINN, and PINN-$k$ on the same held-out rows. This tests whether error clusters at the extremes.
\item Exploratory compartment check: One compartment flagged for an unusually large parameter departure is tracked across all five folds (held out in one fold, included in training for the other four). This test separates a genuine site effect from a fold-specific generalization failure.
\item Yield-class comparison: Each model's personalised parameter deviation is compared with that plot's yield class, a Forest Research site rating excluded from every model's inputs. Tested on the full test set, not only extreme cases, to check whether personalisation tracks this external site-quality measure.
\end{itemize}


\section{Evaluation and robustness} \label{sec:eval-robustness}
This final section explains how model performance and robustness are judged. Model accuracy is evaluated using RMSE, MAE, $R^2$, and Bias. Here $e_i = y_i - \hat y_i$, where $y_i$ is the relevant target: $\bar r_i$ for Q1, $\Delta y_{\max,i}$ for Q2, and top height for Q3.
\begin{equation}
\mathrm{RMSE}=\sqrt{\tfrac{1}{N}\textstyle\sum e_i^2}, \quad
\mathrm{MAE}=\tfrac{1}{N}\textstyle\sum |e_i|, \quad
R^2=1-\tfrac{\sum e_i^2}{\sum (y_i-\bar y)^2}, \quad
\mathrm{Bias}=\tfrac{1}{N}\textstyle\sum e_i.
\label{eq:eval-summary}
\end{equation}
Positive bias indicates under-prediction. We report fold-level means ($\pm\text{SD}$) or pooled metrics across all out-of-fold predictions. For the 4-survey cohort, $R^2$ $95\%$ confidence intervals are calculated using 1,000 compartment bootstrap resamples.

Residual spatial autocorrelation is measured using global Moran's $I$:
\begin{equation}
I=\frac{N}{W}\frac{\sum_i\sum_j w_{ij}(e_i-\bar e)(e_j-\bar e)}{\sum_i(e_i-\bar e)^2}.
\label{eq:morans-i}
\end{equation}
Spatial weights ($w_{ij}$) use binary, row-standardised distance bands matched to each model's own semivariogram range (up to 5,000\,m; 999 permutations). When needed, rows are subsampled to 5,000 to keep this fast to compute. Moran's $I$ is left out if semivariance (the difference between pairs of points at increasing distances) is flat. It is marked provisional if semivariance keeps rising past 5,000\,m. Distances are reported alongside every value. Because each model's band is fitted to its own residual surface, Moran's $I$ values are comparable in direction across models, but not as exact like-for-like magnitudes.

Seed sensitivity checks where ran when computationally feasible rather than across every model. For Elastic Net and XGBoost, we tested eight different spatial fold splits on both the 4-survey and 6-survey cohorts. This tested how sensitive results were to the fold boundaries, rather than to the training process itself. For GNNWR, we re-ran the 6-survey cohort across three different fold splits because GNNWR has no random training seed, while keeping the 4-survey local maps on a single split. For DNN, PINN, and PINN-$k$, we used one seed per fold without repeating splits or training runs; their stability was tested instead using the parameter, physics-weight, and feature tests described above.
For PINN-$k$, we measured how independently the model learns final height ($y_{\max}$) and growth rate ($k$) by calculating their correlation across 2,000 compartment bootstrap samples. We report this alongside $R^2$ for both the main setup and the physics-weight tests, since helping the network tell these two parameters apart is a key goal of adding the physics loss.

%=======================================================================================
%=======================================================================================
\chapter{Results} \label{ch:results}
This chapter answers three questions in order. Q1 links long-term growth shifts to environment and management, checking if different models agree. Q2 tests whether local patterns explain what global models miss. Q3 tests if these signals improve top-height prediction and evaluates a physics-informed network.

\section{Q1: What explains a plot's persistent departure from the shared curve?}
Using the 4-survey cohort, canopy cover and thinning history dominate the attribution in every feature set and every method tested (Table~\ref{tab:results-q1}).
\begin{table}[!htbp]
\centering
\scriptsize
\renewcommand{\arraystretch}{1.2}
\begin{tabular*}{\textwidth}{@{\extracolsep{\fill}}lrrrr@{}}
\toprule
Set & n plots & EN $R^2$ [95\% CI] & XGB $R^2$ [95\% CI] & LMM spatial var. explained \\
\midrule
Set1 (baseline only) & 71{,}766 & 0.220 [0.180, 0.252] & 0.222 [0.183, 0.252] & 0.016$\pm$0.078 \\
Set2 & 71{,}766 & 0.359 [0.310, 0.396] & 0.416 [0.355, 0.461] & 0.204$\pm$0.066 \\
Set3 & 71{,}727 & 0.325 [0.272, 0.365] & 0.375 [0.321, 0.414] & 0.049$\pm$0.048 \\
Set4 & 71{,}330 & 0.358 [0.305, 0.398] & 0.395 [0.338, 0.438] & 0.175$\pm$0.064 \\
\bottomrule
\end{tabular*}
\caption{\footnotesize{Attribution of the persistent mean curve residual (Equation~\ref{eq:mean-cr-residual}),
4-survey cohort, spatial-block split. EN = Elastic Net; XGB = XGBoost; LMM = linear mixed-effects model.
Set1 = canopy cover + thinning only (baseline); Set2 = top-10 ranked environmental predictors; Set3 =
terrain/wind only; Set4 = full environmental set (terrain + climate + soil). $R^2$ 95\% CIs from 1,000
compartment bootstrap resamples. LMM spatial variance explained = share of variance attributable to the
compartment random intercept on top of that set's fixed effects (mean $\pm$ SD across folds, not a
bootstrap CI).}}
\label{tab:results-q1}
\end{table}

These $R^2$ scores are modest: even the best model (Set~2, XGBoost, $R^2=0.416$) leaves over half the growth variance unexplained. The goal is to uncover key environmental drivers rather than make exact plot predictions. Note that $R^2$ is not directly comparable across Q1, Q2, and Q3 because each uses a different target with its own variance. Figure~\ref{fig:q1-fig4-trajectories} illustrates these persistent curve departures in the observed data across six example plots spanning the residual range.

\begin{figure}[!htbp]
  \includegraphics[width=1.1\textwidth]{figures/fig_results/q1/q1_final_fig4_example_trajectories.png}
  \centering
 \caption{\footnotesize{Example plots by residual tier. Observed top height vs.\ age
  (solid) against the pooled Chapman--Richards curve (dashed), for five 4-survey plots: two near
  the 90th percentile of $\bar r_i$ (above curve), one near zero (typical), two near the 10th
  percentile (below curve).}}
\label{fig:q1-fig4-trajectories}
\end{figure}

\subsection{Canopy cover and management dominate, without being a disguised target}
%-----
\textbf{Canopy cover and thinning history dominate attribution across all models and feature sets.}
Canopy cover yields the largest positive Elastic Net coefficient in every configuration (Set~4: $+1.943 \pm 0.034$) and consistently ranks highest in XGBoost feature importance (mean $|\text{SHAP}| = 1.375\text{--}1.442$). Because canopy cover comes from the same airborne LiDAR survey as top height, a key concern is whether it acts as a proxy for the target itself. Three checks evaluate this potential overlap (Figure~\ref{fig:q1-fig1-canopycover-ablation}). First, Pearson correlations between canopy cover and all three targets are moderate ($r = 0.35\text{--}0.47$), showing it is not a direct mathematical duplicate. Second, removing canopy cover from Set~4 reduces $R^2$ by roughly one third for both models (Elastic Net: $0.350 \to 0.231$; XGBoost: $0.388 \to 0.249$; Figure~\ref{fig:q1-fig1-canopycover-ablation}, left), demonstrating that it contains unique structural information not captured by remaining features. Third, candidate replacement variables do not absorb this role; for instance, the importance of $\mathit{dist}_{\text{to\_road}}$ in the XGBoost model drops slightly ($|\text{SHAP}| = 1.003 \to 0.869$; Figure~\ref{fig:q1-fig1-canopycover-ablation}, right) rather than compensating for the excluded signal.

Consequently, canopy cover and thinning history should be interpreted as primary stand-structure baselines. The core attribution question is what true environmental features contribute beyond these structural controls.


\begin{figure}[!htbp]
\centering
 \includegraphics[width=0.8\textwidth]{figures/fig_results/q1/q1_final_fig1_canopycover_ablation.png}
 \caption{\textbf{CanopyCover carries real, non-redundant signal.} Left: pooled $R^2$ (Set4, 4-survey) with vs.\ without CanopyCover for Elastic Net and XGBoost -- both lose roughly a third of their explanatory power. Right: mean $|$SHAP$|$ for CanopyCover, $\mathit{dist\_to\_road}$, precipitation and slope, with vs.\ without CanopyCover -- $\mathit{dist\_to\_road}$'s importance does not rise to compensate, evidence against simple correlation-driven credit transfer. The ``without'' condition is a genuine live re-run, not a transcription of an earlier result.}
\label{fig:q1-fig1-canopycover-ablation}
\end{figure}

%-----
\subsection{Most drivers agree across models; a few genuinely diverge}
\textbf{Most Set~4 predictors agree across models, but six exhibit sharp divergences} anchored by $\mathit{topex}$ (Figure~\ref{fig:q1-fig2-attribution-lenses}; Appendix~\ref{app:nonlinearity-check}). While $\mathit{slope\_degrees}$ and rainfall remain consistently positive across LMM and Elastic Net (slope: $+0.86\text{--}+1.23$), $\mathit{topex}$ splits sharply (LMM $+1.32$ vs.\ Elastic Net $-0.18$ to $+0.06$). This disagreement is not driven by multicollinearity,unstable predictors have lower VIFs (2.1--3.4) than stable rainfall (3.91). It also is not climate-data timing (the static rainfall normal is the stable variable while the survey-matched $\mathit{tas\_mean}$ is unstable. Instead, the divergence reflects spatial confounding: because Elastic Net pools all plots without compartment-level random intercepts, it cannot separate a plot's local shelter (measured by $\mathit{topex}$) from unmeasured regional differences across compartments (such as soil quality or management history), whereas LMM isolates the true within-compartment physical effect. While all unstable predictors exhibit high between-compartment clustering ($\text{ICC} = 0.74\text{--}0.89$), stable rainfall shows an even higher ICC ($0.97$), demonstrating that spatial clustering enables confounding but does not guarantee it in isolation.
%-----
\textbf{Within-compartment variation separates genuine local signals from spatial artifacts.}
Because both models assume linear relationships, their divergence stems from how they handle spatial structure: LMM uses random compartment intercepts, while Elastic Net pools data across the entire landscape. To isolate why coefficients disagree, we calculated within-compartment correlations by mean-centering all predictors and targets within each compartment (subtracting each compartment's mean from its constituent plots). This removes broad regional gradients and tests whether a variable drives growth differences between neighbouring plots.

Across Set~4's 19 predictors, this check sorts the six variables where Elastic Net and the mixed-effects model disagree on sign (Figure~\ref{fig:q1-fig2-attribution-lenses}, left; Appendix~\ref{app:nonlinearity-check}). Only one, $\mathit{topex}$ (green ring), keeps a real local link: its correlation with the target gets \emph{stronger}, not weaker, once compartment differences are removed ($0.08 \to 0.21$). This is why the mixed-effects model's random intercepts can recover a true local effect that Elastic Net, which pools the whole landscape into one line, cannot. The other five (red rings) all show the opposite pattern: their correlation shrinks toward zero once compartment differences are removed -- wind speed ($\mathit{gwa\_weibull\_a\_50m}$, $-0.33 \to -0.14$), temperature ($\mathit{tas\_mean}$, $0.28 \to 0.09$), soil pH ($\mathit{soilgrids\_ph}$, $0.22 \to 0.04$), wetness index ($\mathit{ceh\_twi}$, $-0.11 \to -0.01$), and local relief ($\mathit{local\_relief\_500m}$, $-0.12 \to -0.01$). For these five, the raw correlation mostly reflects which compartment a plot happens to sit in, not a real local effect. Rainfall ($\mathit{chelsa\_bio12\_precip\_mm}$, unringed) is a different case again: Elastic Net and the mixed-effects model already agree on its sign, so it was never one of the six variables this check was built to explain. Predictor instability in the pooled models therefore comes from spatial confounding -- one straight line trying to cover the whole forest -- not from some non-linear pattern the models simply cannot represent.

\begin{figure}[!htbp]
\centering
  \includegraphics[width=\textwidth]{figures/fig_results/q1/q1_final_fig2_attribution_three_lenses.png}
  \caption{\textbf{Three perspectives on environmental attribution (Set~4, 4-survey cohort).} 
\footnotesize \textbf{Left:} Standardized coefficients (mean $\pm$ SD across 5 folds) for Elastic Net (purple) and LMM (grey), with hollow markers for Set~3 and filled markers for Set~4. LMM markers for the six disagreeing variables are ringed by mechanism verdict: green indicates a confirmed local signal ($\mathit{topex}$), while red indicates correlation collapses toward zero once compartment differences are removed. Rainfall ($\mathit{chelsa\_bio12\_precip\_mm}$) is unringed as a stable baseline predictor. \textbf{Right:} XGBoost mean $|\text{SHAP}|$ importance values across the same predictors for comparison.}
\label{fig:q1-fig2-attribution-lenses}
\end{figure}
%-----
\subsection{Spatial clustering falls unevenly, and never fully disappears}
All feature sets leave \textbf{statistically significant spatial clustering in their prediction errors} (distance-band Moran's $I$ at 1385--1739\,m, $p = 0.001$; Figure~\ref{fig:q1-fig3-spatial-clustering}B). Comparing the stand-structure baseline (Set~1) to the full environmental model (Set~4), Elastic Net reduces residual spatial clustering by 29\% ($0.145 \to 0.103$) and XGBoost by 63\% ($0.146 \to 0.054$).

However, XGBoost does not remove spatial clustering smoothly as more data is added. Residual clustering drops sharply at Set~2 ($0.146 \to 0.032$, a 78\% drop), rebounds at Set~3 ($0.077$), and declines again at Set~4 ($0.054$). Ablation tests confirm this bounce is driven by road distance ($\mathit{dist\_to\_road}$) and macroclimate variables ($\mathit{chelsa\_bio12\_precip\_mm}$, $\mathit{tas\_mean}$), which are absent in Set~3. Removing these variables from Set~4 pushes XGBoost's residual Moran's $I$ back up to $0.084$ while costing very little predictive accuracy ($R^2$ drops only from $0.395$ to $0.368$). Remaining residual autocorrelation in both models points to unmeasured environmental drivers or spatially varying processes \cite{gaspard2019Residual_C5_EVAL}.

\begin{figure}[!htbp]
\centering
  \includegraphics[width=\textwidth]{figures/fig_results/q1/q1_final_fig3_spatial_clustering.png}
  \caption{\small{Spatial clustering of model residuals across feature sets (Set~4, 4-survey cohort). 
 Panel A (top): Set~4 Elastic Net residuals averaged by compartment, highlighting leftover spatial error patterns investigated in Q2. Panel B (bottom): Residual Moran's $I$ across feature sets. XGBoost achieves a larger overall reduction than Elastic Net, but the path is non-monotonic, dipping sharply at Set~2 before partially rebounding at Set~3.}}
\label{fig:q1-fig3-spatial-clustering}
\end{figure}

Set~3 (terrain-only) remains the weakest configuration. A category-level refit on Set~4 shows that soil properties add only a small gain ($+0.022$ $R^2$), while macroclimate variables actively reduce test accuracy ($-0.056$ $R^2$; Appendix~\ref{app:category-ablation}). Terrain and edge-effect features drive most of Set~4's predictive performance instead. This challenges the common expectation that forest yield modelling requires broad biophysical climate layers \cite{bontemps2014Predictive_C2_BG}; % [5] Bontemps & Bouriaud (2014) 
in this setup, regional climate grids introduce noise rather than predictive value.

While every environmental model clearly outperforms the stand-structure baseline (Set~1), the 95\% confidence intervals for XGBoost on Set~2 and Set~4 overlap substantially ($[0.355, 0.461]$ vs.\ $[0.338, 0.438]$). As a result, we cannot conclude with confidence that the full environmental suite (Set~4) provides a real operational advantage over the simpler Set~2.


%-----

\subsection*{Q1: Conclusion}
Stand structure and management account for most of a plot's deviation from baseline growth. Among abiotic predictors, slope and precipitation remain consistently stable across methods, whereas six of Set~4's 19 variables show sharp sign disagreements between Elastic Net and the mixed-effects model. This instability stems from spatial pooling: LMM recovers genuine local physical effects for two variables ($\mathit{topex}$ and wind exposure) that Elastic Net obscures, while the remaining four are regional artifacts that collapse once compartment differences are removed. Across feature categories, soil properties add minimal predictive gain ($+0.022$ $R^2$), while macroclimate grids actively degrade test accuracy ($-0.056$ $R^2$). Finally, significant spatial autocorrelation persists across all models' residuals (Moran's $I$ drops 29\% in Elastic Net vs.\ 63\% in XGBoost). Global environmental predictors fail to fully eliminate spatial error patterns, directly motivating the spatially varying coefficients evaluated in Q2.
%=======================================================================================
\section*{Q2: Spatial attribution}
Q1 showed that global models leave strong spatial clustering in their prediction errors. Q2 tests whether a spatially varying model (GNNWR) can capture what these global models miss.

In terms of mean accuracy, GNNWR leads across every feature set (Table~\ref{tab:results-q2}). On Set~4, it reaches $R^2 = 0.294$ ($[0.236, 0.347]$), ahead of Elastic Net ($0.240$, $[0.192, 0.286]$) and XGBoost ($0.250$, $[0.201, 0.295]$). However, broad overlap in confidence intervals means this lead is not clear-cut, and even its best score ($R^2 = 0.320$ on Set~2) leaves most growth variance unexplained. Looking into GNNWR's small advantage shows that because all three models rely almost entirely on canopy cover, GNNWR gains its edge simply by locally smoothing this single main feature, rather than finding subtle environmental trends missed by global baselines.

\begin{table}[h]
\centering
\footnotesize
\renewcommand{\arraystretch}{1.2}
\begin{tabular*}{\textwidth}{@{\extracolsep{\fill}}lrrr@{}}
\toprule
Method & Full Set4 & No CanopyCover & CanopyCover only \\
\midrule
Elastic Net & 0.240 (0.223$\pm$0.100) & 0.042 (0.024$\pm$0.077) & 0.207 (0.190$\pm$0.084) \\
XGBoost & 0.250 (0.231$\pm$0.119) & 0.061 (0.040$\pm$0.094) & 0.210 (0.193$\pm$0.081) \\
GNNWR & 0.294 (0.277$\pm$0.068) & 0.049 (0.026$\pm$0.067) & 0.249 (0.235$\pm$0.030) \\
\bottomrule
\end{tabular*}
\caption{\small{Impact of canopy cover on model performance (Set~4, 4-survey cohort): Pooled $R^2$ (5-fold mean $\pm$ SD in parentheses). All three models collapse near zero without canopy cover ($R^2 < 0.07$). Canopy cover alone accounts for $\sim$80--83\% of each model's predictive power across both global baselines (Elastic Net, XGBoost) and the spatially-varying model (GNNWR).}}
\label{tab:results-q2-ablation}
\end{table}

\subsection{Discussion}
\paragraph*{A noisy target explains the weak signal elsewhere}
Part of the reason other variables barely register is that Q2's target contains heavy noise. The fitting step fixes each plot's growth rate to its assigned yield class and estimates only the maximum height. When a plot grows faster or slower than assumed, that difference does not vanish; it spills straight into the estimated ceiling.
\begin{figure}[H]
\centering
\includegraphics[width=\textwidth]{figures/fig_results/q2/q2_example_plot_curve.png}
\caption{\small{Example plot-level asymptote gap ($\Delta y_{\max}$) and XGBoost attribution. Left: Observed trajectory for Plot 309770, its fitted CR curve, and the yield-class reference curve, with the shaded region showing the observed gap across survey ages 65--81. Right: Local SHAP feature contributions for the same plot, illustrating how individual predictors account for the estimated ceiling departure.}}
\label{fig:q2-example-plot}
\end{figure}

Figure~\ref{fig:q2-example-plot} shows this mechanism for Plot 309770, chosen because it is unusually well-explained rather than typical. Its fitted ceiling reaches 43.0\,m, well above its reference value of 32.8\,m (a $+10.2$\,m gap); XGBoost predicts $+12.4$\,m. On this plot, canopy cover actually lowers the predicted ceiling, in contrast to its positive effect across the broader forest. A typical plot, picked by median prediction error (Appendix Figure~\ref{fig:q2-example-plot-typical}), tells a different story: the model captures only about half the observed gap, with no single factor dominating. This clean example is therefore an exception. Because $\Delta y_{\max}$ merges ceiling error and growth-rate mismatch into one value, only very strong signals—like canopy cover—stand out clearly. Weaker, real environmental effects get lost in this mix, which explains why most other site variables show little impact.

\paragraph*{Q1 and Q2 agree on which variables to trust}
Q1 asks whether Elastic Net and the mixed-effects model (LMM) agree on the direction of a variable's effect (positive or negative). Q2 asks whether GNNWR's local coefficient points the same direction in every forest compartment, or flips between compartments.(Table~\ref{tab:results-q2-reliability} appendix?) Canopy cover and slope pass both checks cleanly: Elastic Net and LMM agree on their direction in Q1, and in Q2 their GNNWR coefficients point the same way in almost every compartment (canopy cover 100\%, slope 97\% of the 231 compartments). Rainfall agrees on direction in Q1 too, and is mostly one-directional in Q2 (87\% of compartments), though less uniformly so than canopy cover or slope. Soil pH fails both checks: Elastic Net and LMM disagree on its direction in Q1, and in Q2 its GNNWR coefficient is positive in barely half the compartments (51\%) and negative in the rest, meaning no real, dependable relationship at either scale.
 
Topographic exposure ($\mathit{topex}$) is the interesting middle case. In Q1, Elastic Net and LMM disagree on its direction because $\mathit{topex}$ has a real, fairly strong effect \emph{within} each compartment (correlation $0.214$) that looks much weaker once compartments are pooled together (correlation $0.080$), so a model that only sees the pooled data gets the direction wrong. Q2 shows the same thing from a different angle: GNNWR's local coefficient for $\mathit{topex}$ is positive in 72\% of compartments but negative in the other 28\% (Appendix Figure~\ref{fig:q2-stability-map}). Because Q1 and Q2 use entirely different statistical methods and still tell the same story, this instability looks like a real feature of the landscape, not an artefact of one particular model.

\begin{figure}[!htbp]
\centering
\includegraphics[width=\textwidth]{figures/fig_results/q2/q2_topex_vs_canopycover_stability.png}
\caption{\small{Spatial stability of GNNWR's local coefficients (Set~4, 4-survey cohort). Each map shows the compartment-average local coefficient on the same colour scale: how much $\Delta y_{\max}$ (m) changes per one-unit increase in the predictor. Left: canopy cover's coefficient stays positive in every compartment. Right: topographic exposure's ($\mathit{topex}$) coefficient is positive in most compartments but flips negative in over a quarter of them. This is why Elastic Net and LMM (Q1) disagree on $\mathit{topex}$'s direction: a global model forced to pick one sign for the whole landscape cannot capture this compartment-by-compartment flip.}}
\label{fig:q2-stability-map}
\end{figure}

A few minor limitations are worth noting. Some terrain features (like wind exposure and slope) show nonlinear patterns—rising before levelling off—which GNNWR's locally linear setup cannot fully capture. However, because canopy cover drives 80--90\% of explained variance, this unmodelled nonlinearity has little practical impact. A tiny fraction of stands (under 1\% of the 4-survey cohort; 266 of 56,526 plots) show extreme trade-offs between fitted growth rates and maximum heights (Appendix~\ref{app:flagged-plot-split}). Finally, testing on the smaller 6-survey cohort produced near-zero accuracy ($R^2 \approx 0$ across all models) due to the small spatial sample size (47 compartments versus 231 in the main cohort).

\subsection*{Q2: Conclusion}
Q2's result is a useful negative finding: letting the model vary its coefficients across space (GNNWR) did not reveal any hidden environmental driver beyond canopy cover, which the simpler global models had already found on their own. The diagnostics explain why. The target itself mixes two different kinds of error into one number, how fast a tree grows, and how tall it eventually gets, which makes any clean spatial pattern harder to find. GNNWR can only fit straight-line relationships within each local area, so it cannot capture terrain effects that curve or level off. In short, these are limits of the method and the target, not evidence that spatial ecology does not matter here.
%=======================================================================================
\section{Q3: Prediction and physics guidance}
Q1 and Q2 found a real environmental signal in growth deviations, Q3 asks whether that signal helps direct height prediction, which model architecture performs best, and whether adding Chapman-Richards physics guidance to a neural network helps further.


% --- Table 1: six-model comparison, 4-survey only ---
\begin{table}[!htbp]
\centering
\scriptsize
\setlength{\tabcolsep}{4pt}
\renewcommand{\arraystretch}{1.2}
\begin{tabular*}{\textwidth}{@{\extracolsep{\fill}}lrrr@{}}\toprule
Model & RMSE${\downarrow}$ & MAE${\downarrow}$ & $R^2{\uparrow}$ \\
\midrule
Linear & $5.168{\pm}0.361$ & $4.056{\pm}0.343$ & $0.580{\pm}0.026$ \\
RF & $4.794{\pm}0.187$ & $3.611{\pm}0.142$ & $0.638{\pm}0.011$ \\
XGBoost (tuned) & $\mathbf{4.548{\pm}0.137}$ & $\mathbf{3.431{\pm}0.123}$ & $\mathbf{0.674{\pm}0.009}$ \\
DNN & $4.681{\pm}0.230$ & $3.533{\pm}0.201$ & $0.655{\pm}0.016$ \\
PINN$^*$ & $4.842{\pm}0.321$ & $3.689{\pm}0.289$ & $0.631{\pm}0.027$ \\
PINN-$k^*$ & $4.927{\pm}0.302$ & $3.745{\pm}0.271$ & $0.618{\pm}0.023$ \\
\bottomrule
\end{tabular*}
\caption{\footnotesize{Top-height model performance on Set~3, 4-survey cohort, mean $\pm$ SD
across five spatial folds ($N=232{,}292$ row-years, $58{,}073$ plots; 39 plots removed from the
$58{,}112$-plot cohort due to missing terrain/wind features). Bold indicates best performance.
$^*$PINN and PINN-$k$ use a corrected forward pass linking predicted parameters directly to
height predictions, evaluated without hyperparameter tuning. Full 6-survey numbers moved to
Appendix~\ref{app:results-q3-6survey} to save space here.}}
\label{tab:results-q3}
\end{table}

% --- Table 2: PINN/PINN-k across feature sets and hyperparameter checks ---
\begin{table}[!htbp]
\centering
\scriptsize
\begin{tabular*}{\textwidth}{@{\extracolsep{\fill}}lrrr@{}}
\toprule
Condition & PINN $R^2$ & PINN-$k$ $R^2$ \\
\midrule
\multicolumn{3}{l}{\textit{Feature set}} \\
No environment & $0.572\pm0.032$ & $0.575\pm0.033$ \\
Set2 (small list) & $0.627\pm0.023$ & $0.622\pm0.017$ \\
Set3 (medium list, headline) & $0.631\pm0.027$ & $0.618\pm0.023$ \\
Set4 (large list) & $0.618\pm0.029$ & $0.620\pm0.020$ \\
\midrule
\multicolumn{3}{l}{\textit{Hyperparameters}} \\
Main result (Set3, reported settings) & 0.631 & 0.618 \\
Learning-rate check & 0.630 & 0.614 \\
Batch-size check & 0.641 & 0.619 \\
Lambda check (8 settings, range) & 0.623--0.632 & 0.612--0.620 \\
\bottomrule
\end{tabular*}
\caption{\small{PINN and PINN-$k$ test $R^2$ across four environment feature-list sizes, and
across four hyperparameter checks on the headline (Set3) configuration. Feature-set values are
5-fold mean $\pm$ SD; hyperparameter values are single trusted runs or, for the lambda check,
the range across 8 settings.}} % source: RESULTS_TABLE.md #5, #3, #3b, #4
\label{tab:pinn-conditions}
\end{table}

% --- Table 3 (was Table 8): plausibility comparison, plain PINN vs. PINN-k. Now the last of
% the 3 tables kept in the main text (see note above on Tables 3-7 moved to appendix). ---
\begin{table}[!htbp]
\centering
\scriptsize
\begin{tabular*}{\textwidth}{@{\extracolsep{\fill}}lrrr@{}}
\toprule
& Plain PINN $y_{max}$ deviation & PINN-$k$ $y_{max}$ deviation \\
\midrule
Mean & $+2.50$\,m & $-0.45$\,m \\
SD & 7.27\,m & 1.36\,m \\
Max & $+25.79$\,m & $+4.95$\,m \\
Min & $-30.45$\,m & $-13.88$\,m \\
Implausible ($<$5\,m or $>$70\,m) & 82 / 58,073 (0.14\%) & 0 / 58,073 (0\%) \\
\bottomrule
\end{tabular*}
\caption{\small{Plain PINN vs.\ PINN-$k$ $y_{max}$ deviation, pooled across all 5 folds
(58,073 plots). PINN-$k$ produces zero implausible heights on the full pooled set.}}
% source: RESULTS_TABLE.md #13 addendum, pooled 5-fold (2026-08-24).
\label{tab:plausibility-comparison}
\end{table}

\subsection{XGBoost achieves the highest predictive accuracy} across all models ($R^2 = 0.674$), ahead of the standard DNN ($0.655$), PINN ($0.631$), and PINN-$k$ ($0.618$; Table~\ref{tab:results-q3}). It wins every fold on the 6-survey cohort (Appendix~\ref{app:results-q3-6survey}) and trails in only 1 of 5 folds on the 4-survey data by just $0.007$ $R^2$. Because requiring six clean surveys filters out steep and high-elevation stands (48 vs.\ 296 compartments), the rest of this chapter focuses on the broader 4-survey cohort. Four possible causes for the DNN-to-XGBoost gap were tested and ruled out: feature interactions (DNN captures 80\% of the nonlinear gain), training sample size (subsampling compartments did not widen the gap), network size (sweeps showed no gain), and out-of-range terrain (XGBoost degraded more than DNN on extreme slopes, error ratio $1.312$ vs.\ $1.243$). This frames the central question for Q3: whether PINNs justify their lower accuracy through biological safety and interpretability.

\subsection{Environmental drivers are balanced, not dominated by one feature}
We tested whether PINN-$k$'s predictions rely on a single hidden feature—similar to how GNNWR leaned on canopy cover for 80--90\% of its signal in Q2—using permutation importance (10 repeats per feature; Figure~\ref{fig:terrain-importance}). Unlike Q2, feature influence here is balanced: the leading predictor (wind exposure) accounts for only 15.2\% of the total importance, with all eleven terrain and climate variables spanning 4.7\% to 15.2\% (evaluated on fold 0). An example on a representative stand (Plot 119937; Figure~\ref{fig:good-example}) demonstrates that this distributed adjustment improves curve fit over the regional benchmark under typical growing conditions.

% Side by side: both narrow (0.7 width alone), both referenced in the paragraph just above.
\begin{figure}[!htbp]
\centering
\begin{subfigure}[t]{0.58\textwidth}
\centering
\includegraphics[width=\textwidth]{figures/fig_results/q3/q3_pinn_terrain_importance.png}
\caption{}
% BUILT, stable. Source: RESULTS_TABLE.md #6.
\label{fig:terrain-importance}
\end{subfigure}
\hfill
\begin{subfigure}[t]{0.38\textwidth}
\centering
\includegraphics[width=\textwidth]{figures/fig_results/q3/q3_pinn_example_plot_curve.png}
\caption{}
% BUILT 2026-08-24 -- extended to show plain PINN alongside PINN-k, same treatment as 77226.
% Underlying graph itself also made narrower (5x5, was 7x5) and legend moved to bottom-right,
% same date -- see build_good_example_plot_curve() in build_q3_redraft_figures.py.
\label{fig:good-example}
\end{subfigure}
\caption{\small{(a) Share of total permutation importance carried by each terrain/wind feature
in PINN-$k$'s personalised $y_{max}$, fold 0. (b) Plot 119937 (compartment 2040, an ordinary
compartment - rank 91 of 178, yield class 16, close to the population mean): both models'
curves fit the observed points better than the population curve, PINN-$k$ most closely of all.
Not an extreme case, shown because personalisation should work sensibly on a typical plot,
not just the flagged ones.}}
\end{figure}

\subsection{A small environmental feature set captures most of the signal}
We evaluated PINN across four environmental configurations under 5-fold cross-validation: a baseline with no site variables, followed by three progressively larger feature sets (Table~\ref{tab:pinn-conditions}, Figure~\ref{fig:env-sweep}). Adding even a minimal feature set (Set~2) produces an immediate performance jump, outperforming the zero-environment baseline across all five folds for both model variants. However, expanding to larger feature sets yields no further gain: paired fold-by-fold comparisons show that Set~2, Set~3, and Set~4 are statistically indistinguishable, with any minor differences sitting well within fold-level noise. Once site variables are included, PINN-$k$ exhibits tighter fold-to-fold consistency than standard PINN, whereas without environmental data both models drop to an identical baseline ($R^2 \approx 0.57$). In practice, while including environmental data is clearly essential, the specific choice among the three candidate feature sets has minimal effect; Set~3 is retained as the default primarily for consistency across the broader dissertation pipeline.

\begin{figure}[!htbp]
\centering
\includegraphics[width=0.75\textwidth]{figures/fig_results/q3/q3_env_sweep_r2.png}
\caption{\small{PINN and PINN-$k$ accuracy against environmental feature-list size.}}
% BUILT, stable. Source: RESULTS_TABLE.md #5.
\label{fig:env-sweep}
\end{figure}

\subsection{How PINN shifts curve parameters and matches yield classes}
To evaluate site-level adaptations, compartments were selected using a predefined ranking rule rather than visual inspection (Table~\ref{tab:productivity-top5}). Among the 178 compartments with $\geq 50$ pooled held-out plots, Compartment 1129 contained the single most extreme plot departure ($+25.8$\,m; rank 1/178), while Compartment 2021 had both the largest mean shift and the most extreme negative plot departure (rank 1/178 on both; Figure~\ref{fig:ymax-deviation-map}). All compartments referenced across this chapter are indexed on a single master reference map (Figure~\ref{fig:all-compartments-reference}, Appendix~\ref{app:results-q3-figures}).

\begin{figure}[!htbp]
\centering
\includegraphics[width=0.8\textwidth]{figures/fig_results/q3/q3_pinn_ymax_deviation_map.png}
\caption{\textbf{Compartment-level asymptote departures under standard PINN (pooled across 5 folds).} Mean $y_{\max}$ deviation from the regional curve, highlighting flagship compartments 1129 and 2021.}
\label{fig:ymax-deviation-map}
\end{figure}

Extending the ranking to the top three over- and under-productive compartments shows strong spatial consistency between model variants (Figure~\ref{fig:hotspot-summary}, Appendix~\ref{app:results-q3-figures}; Table~\ref{tab:productivity-flagships}). However, comparing standard PINN to PINN-$k$ reveals that the two models route site adjustments through different mathematical parameters (Figure~\ref{fig:ymax-k-deviation-maps}). In PINN-$k$, mean $y_{\max}$ stays tightly constrained landscape-wide ($\pm 3$\,m), while the growth rate ($k$) absorbs the primary spatial signal in an asymmetric, upward direction. Four representative compartments (A--D) illustrate these distinct routing behaviors: Compartment 1045 adjusts mainly through $y_{\max}$, Compartment 2057 adjusts through $k$, Compartment 1070 applies a negative $y_{\max}$ shift with neutral $k$, and Compartment 2021 balances adjustments across both parameters.

\begin{figure}[!htbp]
\centering
\includegraphics[width=1.1\textwidth]{figures/fig_results/q3/q3_pinn_k_pooled_ymax_k_deviation_maps.png}
\caption{\small{PINN-$k$, pooled across all 5 folds (58,073 plots): mean $y_{max}$ deviation
(left) vs.\ mean $k$ deviation (right), by compartment. $y_{max}$ stays close to flat almost
everywhere ($\pm$3\,m); $k$ shows a structured, almost entirely one-sided pattern instead --
the same kind of signal, routed through a different parameter than plain PINN uses, and one
that barely moves downward at all. Four compartments labelled A-D on both panels (legend below
the figure) to show that different compartments route the signal differently, not just to
repeat the single most extreme case four times.}}
\label{fig:ymax-k-deviation-maps}
\end{figure}

To test biological plausibility, model predictions were compared against independent yield-class records withheld during training. The five most under-productive compartments aligned with low yield classes in four of five cases, while the top five over-productive compartments matched high yield classes in three of five cases. Across the full held-out test set, model-derived departures correlate significantly with assigned yield classes (Table~\ref{tab:yieldclass-correlation}, Figure~\ref{fig:yieldclass-scatter}, Appendix~\ref{app:results-q3-figures}), showing that the model captures broad ecological gradients alongside local terrain features.

\subsection{Why some compartments disagree with yield class}
The over-productive cases that disagree with regional yield classes show a clear pattern rather than random noise (Table~\ref{tab:productivity-top5}). The two disagreeing compartments (2057 and 1027) are older stands (60 and 72 years) on wind-exposed terrain, whereas the three agreeing compartments are younger (25--35 years) and sheltered. Wind exposure remains a key physical driver, ranking as the leading terrain input across the model (15.2\% permutation importance; Figure~\ref{fig:terrain-importance}). However, stand age cannot simply be dismissed as an extrapolation failure: across all five pooled folds, PINN variants actually achieve lower relative error at older ages (Table~\ref{tab:age-height-error}), leaving the precise role of stand age an open question. Notably, Compartment 2057 also emerges as PINN-$k$'s highest growth-rate departure (rank 1/178), showing that two distinct model parametrisations identify the same exceptional stand.

\subsection{Biological plausibility differs sharply between PINN variants}
The single-parameter PINN occasionally produces extreme extrapolations, projecting asymptote heights well beyond biological limits for a small number of stands (Figure~\ref{fig:implausible-example}). In contrast, the two-parameter PINN-$k$ maintains realistic trajectories across all $58{,}073$ pooled test plots without exception (Table~\ref{tab:plausibility-comparison}, Figure~\ref{fig:plausibility-distribution}). 

Because the standard PINN adjusts only the height ceiling, strong environmental signals force the asymptote upward to fit early growth rates. PINN-$k$ resolves this by channelling rapid early growth into the rate parameter ($k$), leaving the asymptote stable and biologically realistic. This inflation in standard PINN is not a cross-validation artifact: even when trained directly on the affected compartments across four of five folds, the extreme ceiling estimates persisted rather than shrinking (Table~\ref{tab:compartment1129-crossfold}). While it remains open whether these stands represent exceptional site productivity or historical measurement errors, the alignment with independent yield classes supports genuine site variation.

\begin{figure}[!htbp]
\centering
\begin{subfigure}[t]{0.48\textwidth}
\centering
\includegraphics[width=\textwidth]{figures/fig_results/q3/q3_pinn_implausible_example_plot.png}
\caption{}
\label{fig:implausible-example}
\end{subfigure}
\hfill
\begin{subfigure}[t]{0.48\textwidth}
\centering
\includegraphics[width=\textwidth]{figures/fig_results/q3/q3_plausibility_distribution.png}
\caption{}
\label{fig:plausibility-distribution}
\end{subfigure}
\caption{\small{(a) Plot 77226 (compartment 1129): observed heights against the population
curve, plain PINN's curve, and PINN-$k$'s curve. Plain PINN pushes the ceiling to an implausible
77.8m; PINN-$k$ keeps the ceiling modest (53.0m) but personalises growth rate instead
($k$=0.0222, the dataset's single highest value) -- the two curves cross around age 70,
PINN-$k$ fitting the young observed points better before plain PINN's ceiling pulls ahead.
(b) Distribution of $y_{max}$ deviation from the population curve, plain PINN vs.\ PINN-$k$,
pooled across all 5 folds (58,073 plots each). Same underlying signal, very different spread.}}
\end{figure}

\subsection{Model performance differs at age and height extremes}
Model performance shifts noticeably across stand life stages (Table~\ref{tab:age-height-error}, pooled across 5 folds). In the shortest stands ($0$--$10$\,m), the unconstrained DNN is clearly the most accurate (MAE $2.68$\,m vs.\ $3.37$\,m for plain PINN and $3.15$\,m for PINN-$k$). Here, the rigid curve shapes in both PINN variants struggle to match early juvenile growth. In mature stands, the story flips. For trees older than $75$ years ($n = 751$), the DNN becomes the least accurate model ($6.83$\,m vs.\ $6.40$\,m for plain PINN). Without physical constraints, the DNN drifts at older ages, whereas PINN's fixed ceiling anchors predictions to realistic values. This trade-off carries practical weight for forest inventory. Standard volume equations scale height by a power of $1.48$--$1.69$, meaning a $10$\% error in height compounds into roughly a $15$--$17$\% error in timber volume. Because PINN's errors concentrate in juvenile stands rather than mature ones, this volume penalty lands primarily on young plots where timber value is still low.

\subsection{Balancing predictive accuracy against biological safety}
Comparing standard PINN (final height only) with PINN-$k$ (final height and growth rate) reveals a clear trade-off. Standard PINN is slightly more accurate across all four test conditions (Table~\ref{tab:pinn-conditions}). However, PINN-$k$ offers two real advantages: it captures both maximum height and growth rate, and it prevents impossible height predictions entirely (Table~\ref{tab:plausibility-comparison}). The small drop in accuracy happens because tuning two parameters at once is harder for the network: its shared correction trunk has to work roughly twice as hard across all five folds (Table~\ref{tab:trunk-residual}). Even so, the two parameters move together sensibly ($r = 0.43$; Figure~\ref{fig:ymax-k-scatter} in the Appendix), though not symmetrically: growth-rate adjustments are almost entirely positive, while asymptote shifts skew negative. This matches the one-sided pattern seen across individual compartments (Figure~\ref{fig:ymax-k-deviation-maps}), showing the parameters coordinate rather than chase random noise. Testing different learning rates, batch sizes, and physics weights gave the same outcome (Table~\ref{tab:pinn-conditions}), confirming standard PINN's accuracy edge is not just a fluke of one setup. Notably, switching the physics loss off entirely (lambda$=0$) leaves accuracy essentially unchanged (PINN-$k$ $R^2$ $0.620$ vs.\ $0.618$ at the default lambda$=1$), and accuracy drifts mildly downward once lambda exceeds 1 -- physics guidance shapes \emph{how} the model fits, not how well.

\subsection*{Q3: Conclusion}
XGBoost is the most accurate model for predicting top height, but physics-guided networks trade slight accuracy for safer predictions: PINN-$k$ eliminates biologically impossible heights entirely, while plain PINN produces rare but extreme errors (82 of 58,073 test plots, 0.14\%; Table~\ref{tab:plausibility-comparison}). Unlike Q2, Q3's environmental drivers are balanced across terrain and climate features rather than dominated by one variable. Performance also shifts by stand age: the unconstrained DNN leads in young stands, but plain PINN's fixed height ceiling makes it more reliable in mature stands. Therefore, physics guidance shapes \emph{how} the model fits, not how well: it does not boost raw accuracy, but it ensures biologically plausible predictions for the large majority of plots, 99.86\% (57,991 of 58,073) even for plain PINN, and 100\% for PINN-$k$, at a small cost to typical accuracy.

%=======================================================================================




%=======================================================================================
%=======================================================================================
\chapter{Conclusion and Future Work}
\label{ch:conc}

Canopy cover was the strongest correlate of departure from the shared growth curve (Q1). Management variables added smaller signals. Among the environmental variables, only $\mathit{topex}$ and wind exposure showed repeatable local associations after compartment-level differences were considered. These are observational associations, not causal effects. Canopy cover should therefore not be interpreted as a proven management lever.

For the related $\Delta y_{\max}$ target, GNNWR provided only a small improvement over the global models (Q2). Canopy cover accounted for approximately 80--90\% of its explained signal, based on the reduction in $R^2$ when it was removed. The added spatial complexity therefore revealed little further environmental structure.

XGBoost produced the most accurate top-height predictions (Q3). PINN-$k$ had an RMSE approximately 0.4\,m higher, but produced no predictions outside the predefined 5--70\,m plausibility range. Plain PINN produced a small number of large violations. Physics guidance therefore improved plausibility rather than predictive accuracy. This finding applies only within the ages and forest conditions evaluated.

Together, the results address one part of the Forest Digital Twin concept: combining repeated observations and models to explain and predict forest change \cite{baskent2026Review_C1_SS,buonocore2022Proposal_C1_SS}. Q1 and Q2 identify conditions associated with unusual growth. Q3 tests whether those signals support reliable height prediction. Real-time updating, disturbance simulation, management scenarios, and species beyond Sitka spruce remain outside the scope of this work.

\section{Limitations}

\paragraph*{Data and targets}

The environmental layers were coarse relative to individual plots. SoilGrids has a 250\,m resolution, while CHELSA and HadUK climate data are approximately 1\,km \cite{poggio2021SoilGrids_C3_LID,karger2017CHELSA_C3_LID,hollis2019HadUKGrid_C3_LID}. These layers cannot represent fine-scale soil variation or short events such as storms. This may explain the weak and unstable soil and climate signals.

Historical uncertainty also affects the compartment records. Current species, yield-class, and thinning records were applied to earlier surveys and may not describe each compartment accurately at every date. Compartments 2057 and 1027, already flagged in Q3 for disagreeing with their assigned yield class, are exactly the kind of case this could produce, though this project cannot confirm whether their yield-class labels have actually changed over time. Q2's $\Delta y_{\max}$ target has a further limitation: with only four to six observations per plot, differences in growth rate and final height cannot be separated reliably. It is therefore a summary of trajectory departure, not a direct measurement of maximum height.

The six-survey cohort contains only 47 compartments, compared with 231 in the main cohort. Its weak performance may reflect limited spatial variation and unstable holdout folds rather than a fundamental model failure.

\paragraph*{Models and evaluation}

Q1 used one five-fold spatial partition. Its findings, including the slope/$\mathit{topex}$ stability contrast at the centre of Q1's Elastic Net/LMM disagreement, may therefore depend partly on the selected split. GNNWR was tested with additional partitions, but training randomness could not be separated from changes in fold membership, so its reported six-survey instability cannot be cleanly attributed to the model itself. Spatial blocking reduces leakage between nearby observations, but reliable evaluation still depends on having enough independent spatial groups \cite{roberts2017Cross_C5_EVAL,ploton2020Spatial_C5_EVAL}.

Elastic Net and the mixed-effects model assume linear relationships unless nonlinear terms are supplied. This assumption does not fully match the curved responses observed for slope and wind exposure, which qualifies Q1's finding that slope's coefficient sign agrees across models -- agreement on sign does not rule out both models missing the same curvature, though this is likely a minor effect here since canopy cover dominates the explained variance. GNNWR was not systematically tuned, so its small advantage over the global models in Q2 remains uncertain. PINN and PINN-$k$ were also evaluated on one cohort and one training run. The lambda ablation behind Q3's headline claim that physics guidance shapes fit rather than accuracy has not been checked on the six-survey cohort or across seeds, and the separate contribution of the Chapman--Richards branch was not isolated.

\section{Future Work}

The first priority is stronger evaluation. Future experiments should repeat all models across several seeds, tune GNNWR and PINN-$k$ (each GNNWR run took roughly two hours, so cheaper equivalent models with a similar locally-varying structure are also worth testing), and test whether the Chapman--Richards branch contributes beyond the unconstrained neural component. Training on earlier surveys and testing on later surveys would measure temporal prediction -- most naturally on the six-survey cohort, since it has the most years per plot, though that is also the cohort with too few compartments (47) for reliable spatial evaluation, so a temporal split would inherit the same small-sample problem -- while evaluation in another forest, and in other commercially important species beyond Sitka spruce, would test geographic and species transfer.

Future attribution should also examine different forms of growth departure. Separate targets could represent persistently high or low growth, early acceleration, later plateauing, disturbance, and recovery. More surveys would help distinguish changes in growth rate from changes in height potential.

A fuller Forest Digital Twin could move from plot summaries to individual trees. Multi-temporal LiDAR can track tree-level growth and loss, although matching trees between surveys remains difficult \cite{duncanson2018Monitoring_C3_LID}. Using complete three-dimensional point clouds could represent crown size, competition, structural damage, and mortality rather than top height alone \cite{weiser2022Individual_C3_LID}.

Finally, finer soil and climate measurements could be combined with satellite embeddings such as AlphaEarth. These embeddings integrate spatial, temporal, and remote-sensing information, but would reduce direct interpretability and still require spatial and temporal validation \cite{brown2025AlphaEarth_C4_NN}.

%=======================================================================================
%=======================================================================================
\bibliographystyle{plain}
\bibliography{mybibfile}
%=======================================================================================

\appendix
\chapter{Appendix}  \label{chap:appendix}
% \TODO: create appendix secions and labels for each section.

\begin{figure}[!htbp]
\centering
\includegraphics[width=0.65\textwidth]{figures/fig_results/q3/q3_pinn_ymax_k_scatter.png}
\caption{\textbf{Covariance between personalised parameters in PINN-$k$ (pooled across 5 folds).}
Plot-level asymptote departures ($\Delta y_{\max}$) plotted against growth-rate departures ($\Delta k$) relative to each fold's regional baseline ($r = 0.43$). The positive relationship demonstrates coordinated biological scaling across diverse stand conditions.}
\label{fig:ymax-k-scatter}
\end{figure}



%=======================================================================================
\end{document}
%=======================================================================================

